\chapter{近轴成像、渐变介质与光程}
\section{符号与单球面光焦度}\label{th:o1power}
取光沿$z$正向传播，高度$y$向上，近轴倾角$\theta\simeq\dd y/\dd z$有正负。球面顶点为坐标原点，曲率中心在$+z$侧时$R>0$，否则$R<0$。物方、像方折射率为$n_1,n_2$，定义球面光焦度
\[
 \Phi=\frac{n_2-n_1}{R}.
\]
球面高度$y$处法线倾角近似$-y/R$，斯涅尔定律线性化给$n_1(\theta_1+y/R)=n_2(\theta_2+y/R)$，所以
\begin{equation}\label{eq:o1surface}
 n_2\theta_2=n_1\theta_1-\Phi y.
\end{equation}
若另把朝物方的几何半径记为$r>0$，该面$R=-r$，同式就写为$\Phi=(n_1-n_2)/r$。只要曲率和角度约定一致，两种写法给出同一折射方向。

取实物距$u>0$为物点到顶点沿传播方向的距离，实像距$v>0$为顶点到像点的距离，虚物、虚像用负值。轴上物点发出的近轴光有$\theta_1\simeq y/u$，会聚到轴上像点需$\theta_2\simeq-y/v$。代入式\eqref{eq:o1surface}得
\begin{equation}\label{eq:o1gauss}
 \frac{n_1}u+\frac{n_2}v=\Phi.
\end{equation}
前焦距$f_a=n_1/\Phi$，后焦距$f_b=n_2/\Phi$。它们一般不等，二者符号都由$\Phi$确定。将$n_1=\Phi f_a$、$n_2=\Phi f_b$代入成像式并乘$uv/\Phi$，得$f_av+f_bu=uv$，即$(u-f_a)(v-f_b)=f_af_b$。令$x_u=u-f_a$、$x_v=v-f_b$，就有牛顿式$x_ux_v=f_af_b$。这里$x_u,x_v$是相对焦点的偏移，已不是顶点物像距。

\Diagram{fig:o1surface}{球面折射的近轴光线和有向曲率}{\OptAxis\draw[curve](3,-1.6)arc(205:155:3.785);\fill(0,0)circle(1.3pt)node[below]{物点};\fill(6.4,0)circle(1.3pt)node[below]{$C$};\draw[guide](2.7,1)--(6.4,0);\draw[->](0,0)--(2.7,1);\draw[->](2.7,1)--(5.5,.7);\node at(1.2,1.6){$n_1$};\node at(4.8,1.6){$n_2$};}

\begin{example}\label{ex:o1surface}
折射率$n_1=1$的空气中有轴上点光源，位于顶点左方$u=30\,\mathrm{cm}$。右侧玻璃折射率$n_2=1.5$，界面曲率中心在右侧，$R=10\,\mathrm{cm}$。求像距、焦距，并求物距改为10 cm时的像。
\Diagram{fig:o1surfaceq}{空气射入凸球面，曲率中心在玻璃侧}{\OptAxis\draw[curve](3,-1.4)..controls(2.3,-.5)and(2.3,.5)..(3,1.4);\fill(0,0)circle(1.4pt)node[below]{$S$};\fill(4,0)circle(1.4pt)node[below]{$C$};\node at(1.2,1.2){空气};\node at(4.9,1.2){玻璃};}
\end{example}
\begin{solution}
采用图\ref{fig:o1surfaceq}的有向半径：曲率中心在右，故$R=+10\,\mathrm{cm}$。球面光焦度及前后主焦距大小为
\[
 \Phi=\frac{n_2-n_1}R=.05\,\mathrm{cm}^{-1},\quad
 f_a=\frac{n_1}\Phi=20\,\mathrm{cm},\quad
 f_b=\frac{n_2}\Phi=30\,\mathrm{cm}.
\]
其中前焦点在顶点左侧20 cm，后焦点在右侧30 cm，两侧折射率不同，两个焦距大小一般不等。物距为30 cm时，
\[
 \frac{n_2}v=\Phi-\frac{n_1}u
 =\frac1{20}-\frac1{30}=\frac1{60}\,\mathrm{cm}^{-1},
\]
所以$v=90\,\mathrm{cm}>0$，折射光在玻璃中会聚。
\Diagram{fig:o1surfacea}{物点移至前焦点内，折射后仍发散，反向延长得到虚像}{\OptAxis\draw[curve](3,-1.5)..controls(2.4,-.5)and(2.4,.5)..(3,1.5);\fill(1.5,0)circle(1.4pt)node[below]{$S$};\draw[->](1.5,0)--(2.55,.8);\draw[->](2.55,.8)--(5.8,1.65);\draw[guide](2.55,.8)--(-.5,0);\node[below]at(-.5,0){虚像方向};}
物距改为10 cm时，同一成像方程给$1.5/v=.05-.1=-.05\,\mathrm{cm}^{-1}$，故$v=-30\,\mathrm{cm}$。图\ref{fig:o1surfacea}中的虚像由折射光反向延长而成，其坐标在物侧，但整束折射后光线仍属于$n_2$一侧的传播解，成像方程仍使用像方折射率$n_2$。

可用任意近轴高度$y$核对：入射倾角近似为$\theta_1=y/u$，出射约化倾角为$n_2\theta_2=n_1y/u-\Phi y=-n_2y/v$。对实像$v>0$，出射角与$y$反号，向轴靠拢；对虚像$v<0$，出射角与$y$同号，向外发散。这个符号判断直接连接图中的光线走向与成像方程。
\end{solution}

\section{约化光线矩阵与复合系统}\label{th:o1matrix}
用列向量$\vct q=(y,n\theta)^{\mathsf T}$表示近轴光线。厚度$d$、折射率$n$的均匀层只改变高度；球面只改变约化倾角，所以
\[
 T(d,n)=\begin{pmatrix}1&d/n\\0&1\end{pmatrix},\quad
 S(\Phi)=\begin{pmatrix}1&0\\-\Phi&1\end{pmatrix}.
\]
光先经历的矩阵写在最右侧。各矩阵行列式均为1。两折射面相隔$d$、中间折射率$n_3$时，
\[
 M=S(\Phi_2)T(d,n_3)S(\Phi_1)
 =\begin{pmatrix}1-d\Phi_1/n_3&d/n_3\\-\Phi&1-d\Phi_2/n_3\end{pmatrix},
\]
其中$\Phi=\Phi_1+\Phi_2-d\Phi_1\Phi_2/n_3$。厚度项来自第一面改变倾角，进而改变第二面入射高度；简单相加光焦度只在$d=0$或该项可忽略时成立。

\subsection{主平面的位置}
一般近轴系统写成$M=\bigl(\begin{smallmatrix}A&B\\C&D\end{smallmatrix}\bigr)$。若$C\ne0$，等效光焦度为$\Phi=-C$。设第一主面在第一顶点右方$h_1$，第二主面在最后顶点右方$h_2$，则
\[
 M=T(-h_2,n_2)S(\Phi)T(h_1,n_1).
\]
比较$A,D$得$h_1=n_1(1-D)/\Phi$、$h_2=-n_2(1-A)/\Phi$。主面物像距满足$n_1/u_H+n_2/v_H=\Phi$，而不是从镜筒边缘量距离。

若$C=0$，系统无焦，不能用无穷大焦距加上发散的主面位置来计算实际光线。此时应直接使用有限的$A,B,D$矩阵，研究角放大率、束径与光束传播。无焦系统仍由束高、倾角之间的有限矩阵关系确定。

\begin{example}\label{ex:o1pair}
空气中两个薄透镜焦距为$f_1,f_2$，间隔$d=f_1+f_2+\Delta$。求等效焦距、两主面位置及其有向间距；讨论$\Delta=0$。
\Diagram{fig:o1pairq}{两薄透镜间隔$d$及各自焦点的间隔$\Delta$}{\OptAxis\draw[<->](1,-1.3)--(1,1.3);\draw[<->](5,-1.3)--(5,1.3);\node[above]at(1,1.3){$f_1$};\node[above]at(5,1.3){$f_2$};\draw[<->](1,-1.8)--(5,-1.8)node[midway,below]{$d$};\fill(2.7,0)circle(1.3pt)node[above]{$F'_1$};\fill(3.8,0)circle(1.3pt)node[above]{$F_2$};}
\end{example}
\begin{solution}
将图\ref{fig:o1pairq}的两透镜作为$\Phi_i=1/f_i$代入，
\[
 \Phi=\frac{f_1+f_2-d}{f_1f_2}=-\frac\Delta{f_1f_2},\quad f=-\frac{f_1f_2}\Delta.
\]
若$\Delta\ne0$，$A=1-d/f_1$、$D=1-d/f_2$，故从第一透镜位置$z=0$计，$z_{H_1}=fd/f_2$、$z_{H_2}=d-fd/f_1$。两主面的有向间距为
\begin{align*}
 z_{H_2}-z_{H_1}
 &=d-fd(1/f_1+1/f_2)\\
 &=d\left[1-\frac{f_1+f_2}{f_1+f_2-d}\right]
 =\frac{d^2}\Delta.
\end{align*}
\Diagram{fig:o1paira}{主面可能在实体透镜外侧；示例取$f_1=f_2=1,d=3$}{\draw[axis](-3.7,0)--(6.7,0)node[right]{$z$};\foreach\x/\t in{0/L_1,3/L_2}{\draw[<->](\x,-.8)--(\x,.8);\node[above]at(\x,.8){$\t$};}\foreach\x/\t in{-3/H_1,6/H_2}{\draw[guide](\x,-1.2)--(\x,1.2);\node[above]at(\x,1.2){$\t$};}}
图\ref{fig:o1paira}取$\Delta=1$，于是$f=-1$、$H_1$在$-3$、$H_2$在6，主面间隔9。主面是等效光学基准，允许在系统外甚至交叉。若$\Delta=0$，$\Phi=0$，矩阵仍有限，形成无焦系统；主面公式分母为零，已失去适用条件。
\end{solution}

\begin{example}\label{ex:o1ball}
空气中有半径$r$、折射率$n>1$的玻璃球。求等效焦距、主面及背焦距。轴上物点位于球心左方$nr$，求其像的位置并讨论$n=3/2$两侧的情况。
\Diagram{fig:o1ballq}{玻璃球的两个顶点与球心}{\OptAxis\draw[curve](3,0)circle(1.3);\fill(3,0)circle(1.2pt)node[below]{$O$};\fill(.2,0)circle(1.4pt)node[below]{$S$};\node[above]at(3,1.3){$n$};\node[below]at(1.7,0){$V_1$};\node[below]at(4.3,0){$V_2$};}
\end{example}
\begin{solution}
图\ref{fig:o1ballq}的第一面$R_1=r$、$n_1=1,n_2=n$，第二面$R_2=-r$、$n_1=n,n_2=1$，两面光焦度都是$(n-1)/r$。厚度$2r$、中间折射率$n$，所以
\[
 \Phi=\frac{2(n-1)}r-\frac{2r}n\frac{(n-1)^2}{r^2}
 =\frac{2(n-1)}{nr},\quad f=\frac{nr}{2(n-1)}.
\]
$A=D=(2-n)/n$，主面位置$h_1=r,h_2=-r$，两主面均经过球心。背焦距从后顶点起算为$f-r=r(2-n)/[2(n-1)]$。
\Diagram{fig:o1balla}{从球心主面量物像距，图示$3/2<n<3$时的球外实像}{\OptAxis\draw[curve](2.8,0)circle(1);\draw[guide](2.8,-1.4)--(2.8,1.5);\fill(0,0)circle(1.4pt)node[below]{$S$};\fill(6.1,0)circle(1.4pt)node[below]{$S'$};\draw[<->](0,-1.8)--(2.8,-1.8)node[midway,below]{$nr$};\draw[<->](2.8,-1.8)--(6.1,-1.8)node[midway,below]{$nr/(2n-3)$};}
图\ref{fig:o1balla}的主面物距$u_H=nr$。由$1/u_H+1/v_H=1/f$，得$v_H=nr/(2n-3)$。主面位于球心，判断出射束的实虚像还须改从后顶点量距：
\[
 v_{V_2}=v_H-r=\frac{(3-n)r}{2n-3}.
\]
当$3/2<n<3$，$v_{V_2}>0$，折射光在球外会聚成实像。$n=3/2$时出射平行、像在无穷远；$1<n<3/2$时像在球心左侧，为虚像。$n=3$时近轴像落在后顶点。若$n>3$，$0<v_H<r$，像虽在球心右侧，却位于后顶点之前，出射束反向延长相交于球内，仍为虚像。主面像距的正负不能单独代替出射束的会聚、发散判断。
\end{solution}

\section{费马原理与力学类比}\label{th:o1analogy}
在静态、各向同性介质中，几何光线使光程$\int n\dd s$取驻值。变分得到
\begin{equation}\label{eq:o1ray}
 \frac{\dd}{\dd s}(n\vct t)=\nabla n,
\end{equation}
其中$\vct t=\dd\vct r/\dd s$是单位切向量。对分层界面切向积分，由平移对称方向的“广义动量”连续，得$n\sin\theta$守恒，正是折射定律。

为了求轨迹，可以引入虚拟力学粒子，势能$U=-mv_0^2n^2/2$、总能量为零。它的虚拟速度为$|\dot{\vct r}|=v_0n$，运动方程$\ddot{\vct r}=v_0^2\nabla(n^2)/2$。零能量轨道的缩减作用量$\int p\dd s=mv_0\int n\dd s$，所以轨迹与光线相同。但虚拟速度与真实光速$c/n$成反比，不能把虚拟时间当传播时间。若虚拟参数为$q$，
\[
 \dd t_{\rm phys}=\frac{n\dd s}{c}=
 \frac{v_0n^2}{c}\dd q.
\]
对于水波等其他波，$n$可定义为参考波速除以局部相速度；真实时间应按当地群速度计算，浅水非色散模型中群速度等于相速度。

\begin{example}\label{ex:o1graded}
介质折射率满足$n^2(y)=n_0^2(1-y^2/a^2)$，只讨论$|y|\ll a$的近轴光线。一束光在$z=0$处高度$y_0$、倾角$\theta_0$。求轨迹和周期；若改为加号，轨迹有什么变化？
\Diagram{fig:o1gradedq}{中心折射率较大，光线弯向轴线}{\OptAxis\draw[guide](0,1.3)--(6,1.3);\draw[guide](0,-1.3)--(6,-1.3);\fill(0,.8)circle(1.4pt)node[left]{$y_0$};\draw[axis](0,.8)--(1,1.1)node[above]{$\theta_0$};\node[above]at(3,1.3){$|y|$增大时$n$减小};}
\end{example}
\begin{solution}
图\ref{fig:o1gradedq}的虚拟参数取$q$，其方程为$\ddot z=0$、$\ddot y=-v_0^2n_0^2y/a^2$。轴向速度为常数$\dot z=v_0n_i\cos\theta_0$，其中$n_i=n(y_0)$。链式法则给$\ddot y=(\dot z)^2y''(z)$，所以有限初角的精确射线方程为
\[
 y''+g^2y=0,\qquad g=\frac{n_0}{a n_i\cos\theta_0}.
\]
方程的通解为$y=A\cos gz+B\sin gz$，其导数为$y'=-gA\sin gz+gB\cos gz$。初值$y(0)=y_0$、$y'(0)=\tan\theta_0$给$A=y_0$、$B=\tan\theta_0/g$。取$|y_0|/a\ll1$、$|\theta_0|\ll1$，有$g\simeq1/a$、$\tan\theta_0\simeq\theta_0$，才得到题设近轴解
\[
 y=y_0\cos(z/a)+a\theta_0\sin(z/a),\qquad\Lambda_z=2\pi a.
\]
\Diagram{fig:o1gradeda}{负二次项产生周期聚焦，正二次项产生离轴增长}{\draw[axis](0,0)--(6.7,0)node[right]{$z/a$};\draw[axis](0,-1.5)--(0,2.5)node[above]{$y$};\draw[curve,domain=0:6.28,samples=100]plot(\x,{cos(\x r)});\draw[guide,domain=0:1.5,samples=40]plot(\x,{cosh(\x)});\node[right]at(1.6,2.3){加号};}
图\ref{fig:o1gradeda}把折射率平方改为加号。$y$方向虚拟加速度随之变号，近轴方程为$y''-y/a^2=0$。特征根为$\pm1/a$，通解为$C\ee^{z/a}+D\ee^{-z/a}$。初值给$C+D=y_0$、$(C-D)/a=\theta_0$，故$C=(y_0+a\theta_0)/2$、$D=(y_0-a\theta_0)/2$，合并为
\[
 y(z)=y_0\cosh(z/a)+a\theta_0\sinh(z/a).
\]
负号分布中折射率由轴向外减小，射线受回复作用；加号分布则把光线推向轴外。对负号分布，近轴振幅为$\sqrt{y_0^2+a^2\theta_0^2}$，须始终远小于$a$，因此不仅要检查初始高度，还要约束初始倾角引起的振幅。对加号分布，即使初始条件近轴，振幅也可能随传播增长而失效；应在实际截面和近轴允许长度内使用双曲解。
\end{solution}

\subsection{径向介质的角动量}\label{th:o1central}
若$n=n(r)$，虚拟力为中心力，$\vct r\times m\dot{\vct r}$守恒。除以$mv_0$得到
\[
 b=n(r)r\sin\psi,
\]
其中$\psi$是光线切向与径向的夹角，选取$\varphi$递增的绕行方向使$b>0$。弧长参数沿传播方向增加，有
\[
 \frac{\dd r}{\dd s}=\cos\psi
 =\sigma\sqrt{1-\frac{b^2}{n^2r^2}},\qquad
 r\frac{\dd\varphi}{\dd s}=\sin\psi=\frac b{nr},
\]
其中向外支$\sigma=+1$、向内支$\sigma=-1$。在$\dd r/\dd s\ne0$处，两式相除得到
\[
 \frac{\dd\varphi}{\dd r}
 =\sigma\frac{b}{r\sqrt{n^2(r)r^2-b^2}}.
\]
若有孤立最近点，则该点$\cos\psi=0$，故$b=n(r_t)r_t$；转点处不能继续用$r$作单调参数，应接续向内、向外两支。这一等式也可能对应圆轨道等退化情形，须代回射线方程判断。$b=0$为径向光线，直接使用$\dd\varphi/\dd s=0$。若根号内为负，则该区域没有此几何射线；波的穿透须由波动方程处理。

\begin{example}\label{ex:o1wake}
浅水区水深$h(y)=4h_0/(5-y^2/y_0^2)$，岸线$y=0$，船沿$y=y_0$向$+x$以$V_b=2\sqrt{gh_0}$运动。船在$t=0$经过$A=(0,y_0)$。由该处向岸传播的稳态船行波射线与岸线法向夹角$\theta_0$满足$\sin\theta_0=c_A/V_b$，$c_A=\sqrt{gh_0}$。求射线首先到岸的位置、传播时间及届时船的位置。
\Diagram{fig:o1wakeq}{船道、水深变化与向岸初始射线}{\draw[axis](-.3,0)--(6,0)node[right]{$x$};\draw[axis](0,-.2)--(0,3.2)node[above]{$y$};\draw[guide](0,2)--(5.5,2);\fill(0,2)circle(1.6pt)node[left]{$A$};\draw[axis](.2,2.4)--(1.7,2.4)node[above]{$V_b$};\draw[axis](0,2)--(.9,.45);\node[left]at(0,1){$y_0$};\node[below]at(3,0){岸线};}
\end{example}
\begin{solution}
局部浅水波速$c_w(y)=\sqrt{gh(y)}$。取相对折射率$n(y)=c_A/c_w(y)=\tfrac12\sqrt{5-y^2/y_0^2}$，在$A$处$n_A=1$。图\ref{fig:o1wakeq}初始$\sin\theta_0=1/2$，所以$n\sin\theta=1/2$守恒，射线向$+x,-y$运动。

虚拟运动满足$\ddot x=0$、$\ddot y=-v_0^2y/(4y_0^2)$，且$\dot x=v_0/2$。消去虚拟时间得$y''(x)=-y/y_0^2$。正弦型通解为$A\cos(x/y_0)+B\sin(x/y_0)$；由$y(0)=y_0$得$A=y_0$，由$y'(0)=B/y_0=-\cot\theta_0=-\sqrt3$得$B=-\sqrt3y_0$，故
\[
 y(x)=y_0\cos(x/y_0)-\sqrt3y_0\sin(x/y_0)
 =2y_0\cos(x/y_0+\pi/3).
\]
首个$y=0$在$x_B=\pi y_0/6$。
\Diagram{fig:o1wakea}{射线先到岸，船继续前进；横轴按$y_0$归一化}{\draw[axis](0,0)--(6.2,0)node[right]{$x/y_0$};\draw[axis](0,-.1)--(0,2.7)node[above]{$y/y_0$};\draw[guide](0,2)--(5.4,2);\draw[curve,domain=0:.5235988,samples=50]plot({2*\x},{4*cos((\x+pi/3) r)});\fill(1.0472,0)circle(1.5pt)node[below]{$B:\pi/6$};\fill(4.8736,2)circle(1.5pt)node[above]{$A':(\pi+\sqrt3)/2$};\node[left]at(0,2){$1$};}
图\ref{fig:o1wakea}的真实传播时间不可用虚拟周期直接代替。由$\dd s/\dd x=1/\sin\theta=2n$、$c_w=c_A/n$，得
\begin{align*}
 T&=\int_0^{x_B}\frac{\dd s}{c_w}
 =\frac2{c_A}\int_0^{x_B}n^2(x)\dd x\\
 &=\frac1{2c_A}\int_0^{\pi y_0/6}
 [5-4\cos^2(x/y_0+\pi/3)]\dd x\\
 &=\frac1{2c_A}\left[3x-y_0\sin(2x/y_0+2\pi/3)\right]_0^{\pi y_0/6}\\
 &=\frac1{2c_A}\left(\frac{\pi y_0}{2}+\frac{\sqrt3y_0}{2}\right)
 =\frac{y_0}{4c_A}(\pi+\sqrt3).
\end{align*}
最后一步用$4\cos^2u=2+2\cos2u$逐项积分。船位置$x_{A'}=V_bT=y_0(\pi+\sqrt3)/2$。图为清楚显示船位置，横纵采用同一缩放因子2；$B$的横坐标与船到岸时的横坐标不能互换。
\end{solution}

\section{复振幅与光程相位}\label{th:o1phase}
固定时间约定为$\ee^{-i\omega t}$。实电场是$\operatorname{Re}[\mathcal E(\vct r)\ee^{-i\omega t}]$，复振幅$\mathcal E=A\ee^{i\phi}$同时存储幅度和空间相位。同频同偏振光相加时，先加复振幅再求模平方：$I\propto|\sum_j\mathcal E_j|^2$。不同频率或不相关随机相位的交叉项，经足够长时间平均后可以消失。

几何光学中沿射线传播产生相位$k_0\int n\dd s$，$k_0=2\pi/\lambda_0$。若两光束起始相位相同，传播部分相位差为$2\pi\Delta\mathcal L/\lambda_0$；反射、焦散或其他元件的相位须另加。薄片正常插入一束空气光路，以厚$h$的折射率$n$介质替换同厚空气，附加光程$(n-1)h$。斜入射时路径和出射点都改变，不能继续把物理厚度直接代入正常入射公式。

\begin{example}\label{ex:o1plate}
两等振幅相干光原本在某点相长，一束中垂直插入折射率$n$、厚度$h$的平行薄片。求该点光强及使其第一次变暗的最小正厚度。忽略反射损失。
\Diagram{fig:o1plateq}{一束光中以薄片替换同厚空气}{\draw[->](0,1)--(5,1);\draw[->](0,0)--(5,0);\draw[curve](2,.6)rectangle(2.7,1.4);\node[above]at(2.35,1.4){$n,h$};\node[right]at(5,.5){合束};}
\end{example}
\begin{solution}
图\ref{fig:o1plateq}中两复振幅可取$A$、$A\ee^{i\delta}$，$\delta=2\pi(n-1)h/\lambda_0$。合振幅$2A\ee^{i\delta/2}\cos(\delta/2)$，故$I=4I_1\cos^2[\pi(n-1)h/\lambda_0]$。
\Diagram{fig:o1platea}{两相量夹角为$\delta$，合相量沿角平分线}{\draw[axis](0,0)--(3,0)node[right]{$A$};\draw[axis](0,0)--(1.5,2.598076)node[above left]{$A\ee^{i\delta}$};\draw[very thick,->](0,0)--(4.5,2.598076)node[right]{$\mathcal E$};\draw[guide](3,0)--(4.5,2.598076)--(1.5,2.598076);\node at(.9,.4){$\delta$};}
图\ref{fig:o1platea}的两相量首次反向时$\delta=\pi$，所以$h_{\min}=\lambda_0/[2(n-1)]$，$n>1$。若两束振幅不等，反向时仍有残余强度，不能称完全变暗。若频率不相同，静态薄片也不能把随时间移动的拍频图样变成永久暗点。
\end{solution}
\section{费马原理的变分推导}
\subsection{驻值不必是最短路程}
静态各向同性介质中的光线使光程$\mathcal L=\int_A^B n(\vct r)\dd s$在固定端点的路径变分下取驻值。高折射率区的单位路程耗时较长，所以几何短路未必是光程短路。费马原理只要求一阶变分为零，允许局部最小、鞍值及退化驻值。

用任意单调参数$q$表示路径，拉格朗日量为$\mathcal G=n(\vct r)|\dot{\vct r}|$。变分导数
\[
 \frac{\partial\mathcal G}{\partial\dot{\vct r}}=n\vct t,
 \qquad \frac{\partial\mathcal G}{\partial\vct r}=|\dot{\vct r}|\nabla n.
\]
欧拉—拉格朗日方程除以$\dd s/\dd q=|\dot{\vct r}|$，得到$\dd(n\vct t)/\dd s=\nabla n$。把左侧展开为$(\dd n/\dd s)\vct t+n\dd\vct t/\dd s$，再以$\dd n/\dd s=\vct t\cdot\nabla n$，得曲率方程
\[
 n\frac{\dd\vct t}{\dd s}
 =\nabla n-(\vct t\cdot\nabla n)\vct t=\nabla_\perp n.
\]
只有折射率梯度的横向部分使射线转弯；沿射线方向的变化只改变相位积累和局部速度。弯曲方向朝横向折射率较大处，这与分层界面逐层折射的极限一致。

\subsection{界面的切向条件}
在两均匀介质间界面选交点$Q$，端点$A,B$固定。光程为$n_1|A-Q|+n_2|Q-B|$。让$Q$沿界面切向移动$\delta\vct q$，一阶变化为$(n_1\vct t_1-n_2\vct t_2)\cdot\delta\vct q$。任意切向变化都取驻值，要求切向分量$n_1\vct t_{1\parallel}=n_2\vct t_{2\parallel}$，于是$n_1\sin i=n_2\sin r$。

\Diagram{fig:o1fermat}{界面交点的切向变分给出折射定律}{\draw[curve](-3,0)--(3,0);\fill(-2,1.6)circle(1.5pt)node[left]{$A$};\fill(2,-1.4)circle(1.5pt)node[right]{$B$};\fill(0,0)circle(1.5pt)node[above]{$Q$};\draw[curve](-2,1.6)--(0,0)--(2,-1.4);\draw[guide](-2,1.6)--(.6,0)--(2,-1.4);\draw[->](0,.3)--(.6,.3)node[midway,above]{$\delta q$};}
图\ref{fig:o1fermat}的驻值路径在真实界面法向上不要求约化动量连续，因为折射率在该处跳变。光沿边界的相位必须连续，给出相同切向条件。反射则可由镜像法或在同一介质中改变法向传播方向得到，入射角等于反射角。

\subsection{全反射和几何模型的边界}
从$n_1>n_2$射向低折射率介质时，切向条件可能要求$\sin r>1$，此时没有实传播方向。几何光学把它描述为全反射，但并未说明低折射率侧的场严格为零。波动边界解会出现虚法向波矢、倏逝场及偏振相关反射相位。这些内容在第四讲推导，说明同一切向相位约束可以继续使用，而实射线概念已有适用边界。

\section{主平面如何连接实际光线}
主平面是在把系统矩阵分解为“物侧平移—等效折射—像侧平移”时选出的两个参考面。在相同物像介质中，一条假想入射光从实际入口沿直线延长到第一主面，再在等效薄系统中改变倾角，最后从第二主面延长回实际出口；这个构造保持系统外光线关系，不声称系统内部光线按该折线传播。

把矩阵乘法完整写出，设$h_1$从第一顶点向右、$h_2$从最后顶点向右，
\[
 T(-h_2,n_2)S(\Phi)T(h_1,n_1)
 =\begin{pmatrix}
 1+h_2\Phi/n_2&h_1/n_1-h_2/n_2+h_1h_2\Phi/(n_1n_2)\\
 -\Phi&1-h_1\Phi/n_1
 \end{pmatrix}.
\]
由$C=-\Phi$先定总光焦度，再由$A,D$定两个主面，$B$因行列式条件自动一致。若$C=0$，这套分解失去有限主面解，但原$M$依然完整；无焦望远系统、扩束器应直接用矩阵分析。

\subsection{焦距和顶点焦距的差别}
系统外平行光与后焦点的关系，从出射参考面的$y_o=Ay_i$、$n_2\theta_o=Cy_i$可得后焦点距该面的距离$v_F=-n_2A/C=n_2A/\Phi$。另一方面，主面后焦距为$n_2/\Phi$，第二主面相对出口位置为$h_2=n_2(A-1)/\Phi$，相加正好也是$n_2A/\Phi$。这证明主面公式与实际外光线相交位置一致。

玻璃球的两主面在球心，主面后焦距为$f=nr/[2(n-1)]$，从后表面量起的背焦距则为$f-r$。当$n=2$，平行近轴光的焦点在后顶点；当$n>2$，背焦距为负，总光焦度$\Phi$仍为正。实虚像应由后表面外光线的会聚或发散判断，并注明距离的参考面。

\section{渐变介质力学类比的参数选择}
由力学零能量条件$|\dot{\vct r}|=v_0n$，切向单位矢量$\vct t=\dot{\vct r}/(v_0n)$。对$n\vct t=\dot{\vct r}/v_0$按弧长求导，利用$\dd/\dd s=(v_0n)^{-1}\dd/\dd q$，得到
\[
 \frac{\dd(n\vct t)}{\dd s}
 =\frac{\ddot{\vct r}}{v_0^2n}.
\]
令它等于$\nabla n$，便得$\ddot{\vct r}=v_0^2n\nabla n=\tfrac12v_0^2\nabla(n^2)$。这一等式证明零能量力学轨迹满足费马射线方程。任意正值$v_0$只改变虚拟参数的快慢，不改变空间轨迹。

真实传播时间中的权重由$\dd t=n\dd s/c$和$\dd s=v_0n\dd q$共同给出，即$\dd t=v_0n^2\dd q/c$。船行波模型的当地波速为$c_A/n$，相同计算须把分母$c$换成$c_A$。有色散时，波包到达时间应按群速度计算，相位轨迹未必就是能量传播轨迹。

\subsection{对称性给出的积分常数}
若$n$不依$z$，虚拟势具有轴向平移对称性，$\dot z$及$n\cos\alpha$守恒，其中$\alpha$相对$z$轴量取。若$n$只依$r$，转动对称性给$nr\sin\psi$守恒，其中$\psi$相对径向量取。使用守恒量时，须同时保留角度的几何定义。

\Diagram{fig:o1symmetries}{平移对称守恒轴向分量，转动对称守恒角动量}{\begin{scope}[xshift=-2.6cm]\draw[axis](0,0)--(2.5,0)node[right]{$z$};\draw[curve](0,.2)..controls(1,1.3)and(1.6,1.4)..(2.2,1);\draw[guide](0,1.7)--(2.3,1.7);\node[below]at(1.2,-.5){$n\cos\alpha$};\end{scope}\begin{scope}[xshift=1.5cm]\draw[curve](0,0)circle(1);\draw[->](0,0)--(1.6,.8);\draw[->](1.6,.8)--(1.9,1.8);\node[below]at(.5,-1.5){$nr\sin\psi$};\end{scope}}
图\ref{fig:o1symmetries}中的守恒量限定方向和可达区域，继而用于积分求轨迹。若下一界面破坏相应对称性，应按界面相位匹配确定新区域的积分常数。

\subsection{同一共轭关系的横向与纵向变化}
物距和像距必须从各自主面量取。设物、像侧折射率分别为$n_O,n_I$，系统光焦度为$\Phi$，则$n_O/u+n_I/v=\Phi$。在系统不变时微分，
\[
 -\frac{n_O}{u^2}\dd u-\frac{n_I}{v^2}\dd v=0,
 \qquad \frac{\dd v}{\dd u}=-\frac{n_Ov^2}{n_Iu^2}.
\]
物点轴向坐标取$z_O=-u$、像点取$z_I=v$，所以同向轴坐标的纵向放大率为$\dd z_I/\dd z_O=n_Ov^2/(n_Iu^2)>0$。由横向放大率$m_y=-n_Ov/(n_Iu)$，得
\[
 m_z=\frac{\dd z_I}{\dd z_O}=\frac{n_I}{n_O}m_y^2.
\]
在相同外部介质中，轴向微小长度的放大率是横向放大率的平方。横向倒像的负号不使纵向关系变负；物点向右微移，实像在相应共轭分支上也向右移。
\Diagram{fig:o1longmagnify}{两个相近轴上物点与其共轭像点；箭头表示微小轴向位移}{\draw[axis](-3.3,0)--(3.7,0)node[right]{$z$};\draw[thick](0,-1.2)--(0,1.2)node[above]{系统};\fill(-2.8,0)circle(1.5pt)node[below]{$O$};\fill(-2.3,0)circle(1.5pt)node[above]{$O+\dd z_O$};\fill(2,0)circle(1.5pt)node[below]{$I$};\fill(3,0)circle(1.5pt)node[above]{$I+\dd z_I$};\draw[->](-2.8,-.6)--(-2.3,-.6);\draw[->](2,-.6)--(3,-.6);}
图\ref{fig:o1longmagnify}只表示微分共轭，不把有限厚度物体的所有截面当作具有同一个放大率。对有限景深，应分别求前后物面的像距；当物面跨越焦面时，实虚像分支还会改变。矩阵成像式能直接处理这些情况，而只报一个等效焦距不足以定位物像。

当$|m_y|$较大时，物面的小定位误差会被放大为像面位移。几何离焦产生的弥散直径还依赖像方孔径角；衍射斑则由有限孔径决定。用纵向放大率估算焦面位置时，应保持近轴、小位移和同一成像分支三个条件。

\section{成像矩阵与横向放大率}
设物面到像面的总矩阵为$M_{OI}=\bigl(\begin{smallmatrix}A&B\\C&D\end{smallmatrix}\bigr)$。同一物点高度$y_O$发出不同倾角的光，像面高度为$y_I=Ay_O+B n_O\theta_O$。所有出射方向都汇于同一像点，要求$B=0$，横向放大率就是$A$。因为行列式为1，此时$AD=1$，角放大与横向放大有互为倒数的约化关系。

对物像面两条光线，双线性量$y_1n\theta_2-y_2n\theta_1$在每个行列式为1的矩阵下不变，这就是近轴光学不变量。它约束像高与孔径角的乘积：同一系统不可能既把像无代价放大，又保持全部出射角不变。有限孔径、亮度和通光范围需在这个几何约束之外再结合能流研究。

\begin{example}\label{ex:o1matriximage}
近轴系统两端都在空气，总矩阵为$\bigl(\begin{smallmatrix}A&B\\C&D\end{smallmatrix}\bigr)$。物面在入射参考面前$u$，像面在出射参考面后$v$。推导成像条件及放大率，说明它比只给焦距多保留了什么信息。
\Diagram{fig:o1matrixq}{物、像距离分别从系统两参考面量起}{\OptAxis\draw[curve](1.8,-1)rectangle(3.8,1);\node at(2.8,.5){$M$};\draw[->](0,0)--(0,.8);\draw[guide](5.8,-1.2)--(5.8,1.2);\node[below]at(.9,-.5){$u$};\node[below]at(4.8,-.5){$v$};}
\end{example}
\begin{solution}
图\ref{fig:o1matrixq}的光线变量取$(y,p)$，空气中$p=\theta$。由物面先向右传播$u$，再经过系统$M$，最后传播$v$，因此从右向左依次相乘：
\[
 T(v)MT(u)=
 \begin{pmatrix}
 A+vC&Au+B+v(Cu+D)\\ C&Cu+D
 \end{pmatrix}.
\]
同一物点发出的光线具有相同$y_O$、不同$p_O$。要使它们到像面有同一高度，输出高度中$p_O$的系数必须为零，即$Au+B+v(Cu+D)=0$。若$Cu+D\ne0$，
\[
 v=-\frac{Au+B}{Cu+D}.
\]
\Diagram{fig:o1matrixa}{同一物点的两条不同初始倾角光线须落在同一像点}{\OptAxis\draw[curve](2,-1.2)rectangle(3,1.2);\draw[->](0,.8)--(2,1);\draw[->](0,.8)--(2,-.2);\draw[->](3,.5)--(5.5,-.6);\draw[->](3,-.8)--(5.5,-.6);\fill(5.5,-.6)circle(1.4pt);}
图\ref{fig:o1matrixa}中两条入射光线在像面汇于同点，横向放大率为矩阵左上元：
\begin{align*}
 m_y&=A+vC
 =\frac{A(Cu+D)-C(Au+B)}{Cu+D}\\
 &=\frac{AD-BC}{Cu+D}=\frac1{Cu+D}.
\end{align*}
最后使用约化光线矩阵的$\det M=1$。若$Cu+D=0$，行列式又保证$Au+B$不能同时为零，否则矩阵两列线性相关；此时出射为平行束，有限$v$处没有像面。$C=0$的无焦系统仍可用完整矩阵求共轭，无须引入发散的主面坐标。$C$只决定总光焦度，$A,D$还保留焦点相对系统外壳的位置，$B$则影响有限物距的共轭关系。
\end{solution}

\begin{example}\label{ex:o1thicklens}
空气中厚透镜折射率$n$，轴向厚度$d$，前后曲率半径$R_1,R_2$均按本讲有向约定。推导厚透镜制作者公式，并求对称双凸$R_1=R,R_2=-R$的焦距与主面偏移。
\Diagram{fig:o1thickq}{前后曲率中心的方向不同，$R_1>0,R_2<0$}{\OptAxis\draw[curve](2,-1.5)..controls(1.4,-.5)and(1.4,.5)..(2,1.5);\draw[curve](2,-1.5)..controls(3,-.5)and(3,.5)..(2,1.5);\node[above]at(2.2,1.7){$n,d$};}
\end{example}
\begin{solution}
图\ref{fig:o1thickq}中两界面光焦度为$\Phi_1=(n-1)/R_1$、$\Phi_2=(1-n)/R_2$。折射矩阵$S_j=\bigl(\begin{smallmatrix}1&0\\-\Phi_j&1\end{smallmatrix}\bigr)$，玻璃内传播矩阵$T=\bigl(\begin{smallmatrix}1&d/n\\0&1\end{smallmatrix}\bigr)$，故
\[
 S_2TS_1=\begin{pmatrix}
 1-d\Phi_1/n&d/n\\
 -\Phi_1-\Phi_2+d\Phi_1\Phi_2/n&1-d\Phi_2/n
 \end{pmatrix}.
\]
由$C=-\Phi$读出厚透镜制作者公式
\[
 \Phi=(n-1)\left[\frac1{R_1}-\frac1{R_2}
 +\frac{(n-1)d}{nR_1R_2}\right].
\]
\Diagram{fig:o1thicka}{主面向透镜内部移动，两个偏移相反}{\OptAxis\draw[guide](2,-1.5)--(2,1.5);\draw[guide](3,-1.5)--(3,1.5);\draw[very thick](2.2,-1.1)--(2.2,1.1);\draw[very thick](2.8,-1.1)--(2.8,1.1);\node[above]at(2,1.6){$V_1$};\node[above]at(3,1.6){$V_2$};\node[below]at(2.2,-1.1){$H_1$};\node[below]at(2.8,-1.1){$H_2$};}
图\ref{fig:o1thicka}采用$R_1=R$、$R_2=-R$，所以$\Phi_1=\Phi_2=\Phi_s=(n-1)/R$，$\Phi=2\Phi_s-d\Phi_s^2/n$、$f=1/\Phi$。以两个顶点分别为基准的主面有向偏移为
\[
 h_1=\frac{1-D}\Phi=\frac{d\Phi_s}{n\Phi},\qquad
 h_2=\frac{A-1}\Phi=-\frac{d\Phi_s}{n\Phi}.
\]
常规会聚情形$\Phi>0$，前主面向右偏、后主面向左偏；背焦距是从后顶点量到后焦点，故为$f+h_2$，而非$f$。令$d\to0$，两主面均回到薄透镜平面，$\Phi\to2(n-1)/R$。若参数使$\Phi=0$，总矩阵仍可使用，但这些除以$\Phi$的焦距、主面公式不适用，应按无焦系统处理。
\end{solution}

\begin{example}\label{ex:o1apparentdepth}
水中点物体距平界面深$h$，水折射率$n$，空气中观察。求近轴表观深度，并用球面成像的$R\to\infty$极限和射线几何两种方法核对。
\Diagram{fig:o1depthq}{折射光的反向延长给出水中的虚像}{\draw[curve](-2.5,0)--(2.5,0);\fill(0,-2)circle(1.5pt)node[below]{物点};\draw[->](0,-2)--(1,0)--(2.3,1.4);\draw[guide](0,-2)--(0,1.7);\node[left]at(-2.5,-.8){$n$};\node[left]at(-2.5,.8){空气};}
\end{example}
\begin{solution}
沿向上方向取光轴，图\ref{fig:o1depthq}物方$n_1=n$、像方$n_2=1$，平面$\Phi=0$。式\eqref{eq:o1gauss}给$n/h+1/v=0$，所以$v=-h/n$，表观深度为$h/n$。
\Diagram{fig:o1deptha}{近轴时折射角变大，延长线交于较浅位置}{\draw[curve](-2.5,0)--(2.5,0);\fill(0,-2)circle(1.5pt)node[left]{$h$};\fill(0,-1.2)circle(1.5pt)node[left]{$h/n$};\draw[->](0,-2)--(1,0);\draw[->](1,0)--(2.5,1.8);\draw[guide](1,0)--(0,-1.2);}
射线在水中角$i$、空气角$r$，$n\sin i=\sin r$。若界面交点距轴$y$，真实深度$h=y/\tan i$，延长交点深度$h_a=y/\tan r$。近轴$r\simeq ni$，所以$h_a\simeq h/n$。图\ref{fig:o1deptha}画的是近轴虚像；大角度时$h_a=h\tan i/\tan r$依角度变，平面折射并不对所有孔径形成严格同一点像。
\end{solution}

\section{渐变介质的传像与精确光线}
\begin{example}\label{ex:o1grinrelay}
近轴渐变介质满足$y''+g^2y=0$，轴上折射率近似$n_0$。求长度$l$的约化矩阵，分别讨论四分之一节距和半节距的用途。
\Diagram{fig:o1grinq}{渐变杆内的周期光线}{\OptAxis\draw[curve](0,-1.2)rectangle(6,1.2);\draw[curve,domain=0:6,samples=80]plot(\x,{.8*cos(60*\x)});\node[above]at(3,1.2){$y''+g^2y=0$};}
\end{example}
\begin{solution}
由初值解$y(l)=y(0)\cos gl+y'(0)\sin gl/g$、$y'(l)=-gy(0)\sin gl+y'(0)\cos gl$，图\ref{fig:o1grinq}的约化矩阵为
\[
 M(l)=\begin{pmatrix}\cos gl&\sin gl/(n_0g)\\-n_0g\sin gl&\cos gl\end{pmatrix}.
\]
四分之一节距$gl=\pi/2$，输出高度由输入倾角决定、输出倾角由输入高度决定；轴上点物发出的光变为轴向平行光。半节距$gl=\pi$，$M=-I$，输入面被倒立等大地传到输出面。
\Diagram{fig:o1grina}{半节距中，来自同一输入点的不同初始倾角光线在反高处重聚}{\draw[axis](0,0)--(6.6,0);\draw[curve,domain=0:6,samples=80]plot(\x,{.7*cos(30*\x)});\draw[guide,domain=0:6,samples=80]plot(\x,{.7*cos(30*\x)+.6*sin(30*\x)});\fill(0,.7)circle(1.4pt);\fill(6,-.7)circle(1.4pt);\node[below]at(3,-1.1){$gl=\pi$};}
图\ref{fig:o1grina}的两光线初始斜率不同，末高却同为$-y_0$，体现$B=0$成像条件。四分之一节距不能同时又说输出面是输入面的等大像；它交换了位置和倾角的信息。若杆端与空气折射，还需接上两个平面界面的角度变换，本约化变量已自动保持$n\theta$连续。
\end{solution}

\begin{example}\label{ex:o1linearn2}
折射率平方在某区域满足$n^2(y)=n_A^2+2ay$，$a$为常数。光从原点沿$+z$水平进入，求虚拟力学轨迹和真实传播时间，限定在折射率为实的区间。
\Diagram{fig:o1linearq}{折射率平方的梯度使光线弯向高折射率侧}{\OptAxis\draw[axis](0,0)--(0,2.7)node[above]{$y$};\draw[->](0,0)--(1.7,0);\node at(3,1.4){$n^2=n_A^2+2ay$};}
\end{example}
\begin{solution}
图\ref{fig:o1linearq}中，虚拟参数$q$满足$\ddot{\vct r}=(v_0^2/2)\nabla n^2$，故$\ddot z=0$、$\ddot y=av_0^2$。光从原点水平进入，$z(0)=y(0)=0$、$\dot z(0)=n_Av_0$、$\dot y(0)=0$。先积分加速度，$\dot y=\int_0^qav_0^2\dd s=av_0^2q$、$\dot z=n_Av_0$；再积分速度，
\[
 z=n_Av_0q,\qquad y=\tfrac12av_0^2q^2
 =\frac{az^2}{2n_A^2}.
\]
此处的速度是对虚拟参数求导所得；真实时间仍由光程积分给出。沿轨迹$y'=az/n_A^2$，且$n^2=n_A^2+a^2z^2/n_A^2=n_A^2(1+y'^2)$。
\Diagram{fig:o1lineara}{轨迹为抛物线，但真实传播时间不是匀加速粒子时间}{\OptAxis\draw[axis](0,0)--(0,2.9)node[above]{$y$};\draw[curve,domain=0:4.5,samples=70]plot(\x,{.12*\x*\x});\draw[guide](3,0)--(3,1.08);\node[below]at(3,0){$Z$};}
真实传播时间为$\dd t=n\dd s/c$，图\ref{fig:o1lineara}中的弧长满足$\dd s=\sqrt{1+y'^2}\dd z=(n/n_A)\dd z$，因此
\begin{align*}
 t(Z)&=\frac1{cn_A}\int_0^Z n^2(z)\dd z\\
 &=\frac1{cn_A}\left[n_A^2z+\frac{a^2z^3}{3n_A^2}\right]_0^Z
 =\frac{n_AZ}c+\frac{a^2Z^3}{3cn_A^3}.
\end{align*}
几何路程的增加和沿途折射率的增加都已包含在$n\dd s$中。若$a<0$，光线向负$y$偏折，但$a y>0$，沿射线仍是折射率增大；时间修正含$a^2$，保持正值。$a$的量纲为长度倒数，$a^2Z^3/c$为时间。以上轨迹只适用于该分布覆盖的区域，遇到界面后须重新匹配出射方向。
\end{solution}

\begin{example}\label{ex:o1luneburg}
半径$R$的球形渐变透镜在球内$n^2(r)=2-r^2/R^2$，球外$n=1$。证明平行于$x$、冲量参数满足$|b|<R$的入射光线会聚于球面右端$(R,0)$，并求从左侧平面$x=-R$到焦点的总光程。
\Diagram{fig:o1luneq}{球边$n=1$，平行光以不同冲量参数$b$进入}{\draw[axis](-3,0)--(3,0)node[right]{$x$};\draw[axis](0,-2.2)--(0,2.4)node[above]{$y$};\draw[curve](0,0)circle(1.8);\foreach\y in{-.9,0,.9}{\draw[->](-2.8,\y)--({-sqrt(3.24-\y*\y)},\y);}\node[right]at(2,1.4){$n^2=2-r^2/R^2$};}
\end{example}
\begin{solution}
取图\ref{fig:o1luneq}二维截面，入点$\vct r_0=(-a,b)$，$a=\sqrt{R^2-b^2}$。球边折射率连续且等于1，入射方向不跳变。虚拟力为$\ddot{\vct r}=-(v_0^2/R^2)\vct r$，初始速度$v_0\hat{\vct x}$。令$\eta=v_0q/R$，精确解
\[
 x=-a\cos\eta+R\sin\eta,\qquad y=b\cos\eta.
\]
在$0<\eta<\pi/2$，$r^2=R^2-2aR\sin\eta\cos\eta<R^2$，光线确在球内；首个再次到球面时$\eta=\pi/2$，得到$x=R,y=0$，与$b$无关。
\Diagram{fig:o1lunea}{精确光线在球面同一点相交；这里无需近轴条件}{\draw[axis](-2.8,0)--(2.6,0);\draw[curve](0,0)circle(1.8);\foreach\b in{-.9,0,.9}{\pgfmathsetmacro{\a}{sqrt(3.24-\b*\b)}\draw[curve,domain=0:90,samples=70]plot({-\a*cos(\x)+1.8*sin(\x)},{\b*cos(\x)});}\fill(1.8,0)circle(1.7pt)node[right]{$F$};}
沿轨迹$n^2=1+(a/R)\sin2\eta$，真实光程为
\[
 ct=R\int_0^{\pi/2}\left(1+\frac aR\sin2\eta\right)\dd\eta
 =R\left[\eta-\frac{a}{2R}\cos2\eta\right]_0^{\pi/2}
 =\frac{\pi R}{2}+a.
\]
从左平面到入点的空气程为$R-a$，相加得总光程$R(1+\pi/2)$，不依冲量参数。图\ref{fig:o1lunea}的几何聚焦同时满足等光程。严格掠射的$|b|=R$不属于穿入球内的光线，以上结论可取其从内侧逼近的极限。
\end{solution}

\section{光程、相位与短波近似}
设标量场满足$\nabla^2U+k_0^2n^2U=0$，写$U=A\ee^{ik_0S}$，$A,S$缓变且$S$为实。第一次求导给$\nabla U=\ee^{ik_0S}(\nabla A+ik_0A\nabla S)$；再用乘积法则求散度，
\[
 \nabla^2U=\ee^{ik_0S}\left[
 \nabla^2A+2ik_0\nabla A\cdot\nabla S
 +ik_0A\nabla^2S-k_0^2A|\nabla S|^2\right].
\]
代回波动方程并除去指数，得$\nabla^2A+ik_0(2\nabla A\cdot\nabla S+A\nabla^2S)+k_0^2A(n^2-|\nabla S|^2)=0$。短波渐近展开中，先令$k_0^2$阶系数为零，得到程函方程$|\nabla S|^2=n^2$；再令$k_0$阶系数为零，得到输运方程$2\nabla A\cdot\nabla S+A\nabla^2S=0$。后式乘$A$，利用$\nabla(A^2)=2A\nabla A$，可写成$\nabla\cdot(A^2\nabla S)=0$。第一式决定相位面的法线方向和光程，第二式决定光线束收缩时的振幅变化；$\nabla^2A$留在更低阶，焦区等振幅急变区域不能将它忽略。

\begin{example}\label{ex:o1eikonal}
在均匀介质中，分别求平面波和球面出射波的程函$S$及振幅随传播距离的变化，解释焦散附近为何几何光学振幅会失效。
\Diagram{fig:o1eikonalq}{平行光管截面积不变，球面光管面积随$r^2$增大}{\begin{scope}[xshift=-2cm]\draw[->](0,-.6)--(2,-.6);\draw[->](0,.6)--(2,.6);\node[above]at(1,1){平行光管};\end{scope}\draw[->](2,0)--(4.5,1.2);\draw[->](2,0)--(4.5,-1.2);\node[above]at(3.5,1.7){发散光管};}
\end{example}
\begin{solution}
图\ref{fig:o1eikonalq}的平面波可取$S=nz$，$\nabla^2S=0$。输运方程给$\partial_z A=0$，积分为$A(z)-A(z_0)=\int_{z_0}^z0\dd q=0$，故振幅不随$z$衰减。球面波取$S=nr$，径向拉普拉斯为$\nabla^2S=r^{-2}\partial_r(r^2\partial_rS)=r^{-2}\partial_r(nr^2)=2n/r$。输运方程变为$A'=-A/r$，在$r>0$分离变量，
\[
 \int_{A(r_0)}^{A(r)}\frac{\dd a}{a}
 =-\int_{r_0}^{r}\frac{\dd q}{q},\qquad
 \ln\frac{A(r)}{A(r_0)}=-\ln\frac r{r_0}.
\]
因此$A(r)=A(r_0)r_0/r$，强度随$1/r^2$衰减。
\Diagram{fig:o1eikonala}{会聚光管在焦点截面积趋零，短波振幅公式不能无限使用}{\draw[->,thick](-2,1.2)--(0,0)--(2,-1.2);\draw[->,thick](-2,-1.2)--(0,0)--(2,1.2);\draw[guide](-1,-.6)--(-1,.6);\fill(0,0)circle(1.5pt)node[above]{$F$};\node[below]at(0,-1.7){焦区需波动解};}
球面能流乘球面积保持不变。图\ref{fig:o1eikonala}会聚时光管面积趋零，几何输运式形式上给出发散振幅，表明缓变假设在焦区失效。焦区应改用衍射解，得到有限场强及焦点相移。
\end{solution}

\providecommand{\ExpansionDiagram}[3]{\begin{figure}[H]\centering\begin{tikzpicture}#3\end{tikzpicture}\caption{#2}\label{#1}\end{figure}}
\section{有限孔径、光瞳与可成像的视场}
\subsection{成像条件并不保证光线能够通过}
矩阵给出光线的交点，镜框、光阑和透镜边缘则限制哪些光线能够到达该处。求出像面后，还需逐面计算光线高度，与各面的通光口径比较。

取物面光线坐标为$(y,p)^{\mathsf T}$，其中$p=n_O\theta$。从物面到第$j$个孔径平面的矩阵第一行为$(A_j,B_j)$，该孔径半宽为$a_j$，则通过条件是
\[
 |A_jy+B_jp|\leq a_j.
\]
给定物点高度$y$，每个孔径在$p$轴上限定一个区间，通过全部孔径的光线方向由这些区间的交集确定。$B_j=0$时，该孔径只限制物高；$B_j\ne0$时，它限制该物点可用的光线方向。交集缩小会造成渐晕，交集为空则该物点没有光到达像面。采用同一套有向光线坐标时，这一判定也适用于虚物和无焦系统。

对轴上物点$y=0$，最先限制入射角的孔径是孔径光阑。它通过前方光学元件所成的像称为入瞳，通过后方元件所成的像称为出瞳。光瞳可以是虚像，不能把入瞳一律画在真实光阑处。主光线通过孔径光阑中心；边缘光线通过光阑边缘。前者便于追踪视场，后者便于判断数值孔径。

\ExpansionDiagram{fig:dep-o1-interval}{固定物高时，不同孔径对入射方向的限制形成区间交集}{
\draw[->](0,0)--(8,0)node[right]{$p$};
\draw[very thick](1,.7)--(5,.7);\draw(1,.55)--(1,.85);\draw(5,.55)--(5,.85);\node[left]at(.9,.7){孔径1};
\draw[very thick](3,1.5)--(7,1.5);\draw(3,1.35)--(3,1.65);\draw(7,1.35)--(7,1.65);\node[left]at(2.9,1.5){孔径2};
\fill[gray!20](3,-.25)rectangle(5,.25);\node[below]at(4,-.35){通过区间};
\draw[dashed](3,-.25)--(3,1.6);\draw[dashed](5,-.25)--(5,1.6);}
图\ref{fig:dep-o1-interval}中的阴影是两个允许区间的交集。物高变化时，区间通常发生相对平移，所以边缘视场的有效通光量可能低于轴上通光量，即使两个物点都严格成像。

\subsection{光学不变量与通光量的几何限制}
设同一系统中的两条近轴光线为$(y_1,p_1)$、$(y_2,p_2)$。约化矩阵的行列式为1，故
\[
 H=y_1p_2-y_2p_1
\]
在近轴传播和折射中保持不变。直接代入$y_j'=Ay_j+Bp_j$、$p_j'=Cy_j+Dp_j$，展开两乘积，$ACy_1y_2$及$BDp_1p_2$项相消，剩下$H'=(AD-BC)(y_1p_2-y_2p_1)=H$。它是二维光线相空间中的有向面积。选一条从轴上物点出发的边缘光线和一条从轴外物点出发的主光线，在物面得到$|H|=n_O|y_O\theta_O|$，在像面则为$n_I|y_I\theta_I|$。因此
\[
 n_O|y_O\theta_O|=n_I|y_I\theta_I|.
\]
放大像高以后，像方角孔径相应减小。孔径遮挡能删去一部分光线，却不能在保留全部入射光线的同时任意压缩这块相空间面积。这里使用近轴角度；大角度系统应以正弦形式及相应光学条件重新讨论。

\begin{example}两孔径引起的渐晕
空气中薄透镜焦距$f>0$、半口径$a$，物面在透镜前$2f$。透镜后$f$处有半宽同为$a$的光阑。求像面位置、轴上角孔径，以及物高$h$变化时可通过的方向区间和无光视场边界。
\ExpansionDiagram{fig:dep-o1-aperture-q}{透镜与后置光阑分别限制光线高度}{
\draw[->](-3.5,0)--(4,0)node[right]{$z$};
\draw[<->,thick](0,-1)--(0,1);\node[above]at(0,1.2){透镜};\draw[<->](-.45,0)--(-.45,1)node[midway,left]{$a$};
\draw[very thick](1.5,1)--(1.5,1.7);\draw[very thick](1.5,-1)--(1.5,-1.7);\node[above]at(1.5,1.8){光阑};
\draw[->](-3,0)--(-3,.8)node[above]{$h$};\draw[dashed](3,-1.4)--(3,1.4)node[above]{像面};
\draw(-3,-.5)--(0,-.5)node[midway,below]{$2f$};\draw(0,-1.9)--(1.5,-1.9)node[midway,below]{$f$};}
\end{example}
\begin{solution}
由$1/(2f)+1/v=1/f$得$v=2f$，横向放大率为$-1$。图\ref{fig:dep-o1-aperture-q}中像面在光阑之后，光阑本身不改变成像位置。

设物点发射角为$\theta$。到达透镜时高度$y_L=h+2f\theta$，折射后方向为
\[
 \theta_+=\theta-\frac{y_L}{f}=-\frac hf-\theta.
\]
再传播$f$到光阑，高度变为$y_S=y_L+f\theta_+=f\theta$。两个通过条件分别给出
\[
 \frac{-a-h}{2f}\leq\theta\leq\frac{a-h}{2f},
 \qquad -\frac af\leq\theta\leq\frac af.
\]
轴上第一段较窄，所以透镜边缘是孔径光阑，轴上角半宽$a/(2f)$。
\ExpansionDiagram{fig:dep-o1-aperture-a}{方向区间随物高平移；临界时交集收缩为一点}{
\draw[->](0,0)--(7,0)node[right]{$f\theta/a$};
\draw[very thick](1,.5)--(5,.5);\node[left]at(.9,.5){后光阑};
\draw[very thick](2,1.2)--(4,1.2);\node[right]at(4.1,1.2){$h=0$};
\draw[very thick](0,1.9)--(2,1.9);\node[right]at(2.1,1.9){$h=2a$};
\draw[very thick](-1,2.6)--(1,2.6);\node[right]at(1.1,2.6){$h=3a$};
\node[below]at(1,0){$-1$};\node[below]at(3,0){$0$};\node[below]at(5,0){$1$};\draw[dashed](1,.2)--(1,2.8);}
在图\ref{fig:dep-o1-aperture-a}中，透镜允许区间的中心为$-h/(2f)$，半宽$a/(2f)$；后光阑的区间固定。由两个端点比较可得方向区间长度
\[
 \Delta\theta=\begin{cases}
 a/f,& |h|\leq a,\\
 (3a-|h|)/(2f),&a<|h|<3a,\\
 0,&|h|\geq3a.
 \end{cases}
\]
$|h|=3a$时仅剩一条几何边界光线，其连续光束通量为零；$|h|>3a$时完全遮挡。这里求的是一个子午截面的方向宽度，不能直接把它当作三维圆孔系统的光通量比例；三维情况还需积分允许方向的面积。
\end{solution}

\begin{example}像面偏离与弥散斑
仍取焦距$f$、半口径$a$的薄透镜，轴上物点在$2f$处。屏从理想像面沿$+z$移开$\delta$。求几何弥散斑半径；若物点还发生小轴向位移$\Delta z_O$，求一阶合并后的离焦量。忽略衍射，要求$|\delta|,|\Delta z_O|\ll f$。
\ExpansionDiagram{fig:dep-o1-defocus-q}{像点两侧的截面都可能形成弥散斑}{
\draw[->](-.5,0)--(6,0)node[right]{$z$};\draw[<->,thick](0,-1.3)--(0,1.3);
\draw(0,1.2)--(4,0)--(5,-.3);\draw(0,-1.2)--(4,0)--(5,.3);
\fill(4,0)circle(1.4pt)node[below]{$I$};\draw[thick](5,-1)--(5,1)node[above]{屏};\draw[<->](4,.65)--(5,.65)node[midway,above]{$\delta$};}
\end{example}
\begin{solution}
图\ref{fig:dep-o1-defocus-q}中任一到达透镜高度$y_L$的光线，在透镜后为$y(z)=y_L(1-z/(2f))$。屏在$z=2f+\delta$，故$y=-y_L\delta/(2f)$。取$|y_L|\leq a$即得斑半径
\[
 r_g=\frac{a|\delta|}{2f}.
\]
物点坐标向右移动$\Delta z_O$意味着物距由$2f$变为$2f-\Delta z_O$。对$1/u+1/v=1/f$求微分，$-\dd u/u^2-\dd v/v^2=0$，故$\dd v/\dd u=-v^2/u^2$。在$u=v=2f$代入得$\dd v=-\dd u$，所以像点向右移动$\Delta z_I=\Delta z_O$。屏相对新像面的离焦量为$\delta-\Delta z_O$。
\ExpansionDiagram{fig:dep-o1-defocus-a}{物面移动与屏面移动共同决定实际离焦}{
\draw[->](0,0)--(7,0)node[right]{$z$};\fill(2,0)circle(1.4pt)node[below]{$I$};\fill(3,0)circle(1.4pt)node[below]{$I'$};
\draw[thick](5,-1)--(5,1)node[above]{屏};\draw[<->](2,.4)--(3,.4)node[midway,above]{$\Delta z_O$};\draw[<->](3,-.7)--(5,-.7)node[midway,below]{$\delta-\Delta z_O$};}
因此一阶结果为$r_g\simeq a|\delta-\Delta z_O|/(2f)$。图\ref{fig:dep-o1-defocus-a}中两种位移同向且相等时，一阶离焦消失。若需要更高精度，应使用新的像距与新的像方角度；若斑半径已接近衍射尺度，则几何斑不能单独作为实际点像宽度。
\end{solution}

\section{矩阵成像与高斯光束的连接}
\subsection{高斯积分的归一化与带相位形式}
光束传播与相干平均都会用到
\[
 \mathcal G(a,b)=\int_{-\infty}^{\infty}\exp(-ax^2+bx)\dd x,
 \qquad \operatorname{Re}a>0.
\]
先取实数$a>0,b=0$。将积分平方写成二维积分，再取极坐标，
\[
 \mathcal G(a,0)^2=\int_0^{2\pi}\dd\varphi\int_0^\infty
 r\ee^{-ar^2}\dd r
 =2\pi\left[-\frac{\ee^{-ar^2}}{2a}\right]_0^\infty=\frac\pi a.
\]
原积分为正，故$\mathcal G(a,0)=\sqrt{\pi/a}$。对$b$求导并分部积分，
\begin{align*}
 \partial_b\mathcal G
 &=\int_{-\infty}^{\infty}x\ee^{-ax^2+bx}\dd x\\
 &=-\frac1{2a}[\ee^{-ax^2+bx}]_{-\infty}^{\infty}
 +\frac b{2a}\int_{-\infty}^{\infty}\ee^{-ax^2+bx}\dd x
 =\frac b{2a}\mathcal G.
\end{align*}
边界项由$\operatorname{Re}a>0$保证为零。乘积分因子$\exp[-b^2/(4a)]$，有$\partial_b\{\ee^{-b^2/(4a)}\mathcal G\}=0$，以$b=0$的值确定常数，得到
\begin{equation}\label{eq:o1gaussintegral}
 \mathcal G(a,b)=\sqrt{\frac\pi a}\exp\frac{b^2}{4a}.
\end{equation}
对复数$a$取在正实轴上为正的连续平方根。积分与右侧在$\operatorname{Re}a>0$均解析，上式由实正$a$连续延拓到该半平面；它适用于后面含二次相位的收敛积分。

\subsection{光束宽度与波前曲率需要同时传播}
束腰附近的有限光斑会发生衍射扩展，需要同时计算场的振幅与相位。对沿$+z$传播、时间因子为$\ee^{-\mathrm i\omega t}$的标量场，写作$E=U(x,y,z)\ee^{\mathrm ikz}$。在$|\partial_z^2U|\ll|k\partial_zU|$的近轴条件下，Helmholtz方程化为
\[
 2\mathrm ik\partial_zU+\nabla_\perp^2U=0.
\]
均匀介质中试取轴对称解$U=Cq^{-1}\exp[\mathrm ikr^2/(2q)]$。对$z,r$分别求导，
\begin{align*}
 \frac{\partial_zU}{U}&=-\frac{q'}q-\frac{\mathrm ikr^2q'}{2q^2},\\
 \frac{\partial_rU}{U}&=\frac{\mathrm ikr}q,\qquad
 \frac{\partial_r^2U}{U}=\frac{\mathrm ik}q-\frac{k^2r^2}{q^2},\\
 \frac{\nabla_\perp^2U}{U}&=\frac{2\mathrm ik}q-\frac{k^2r^2}{q^2}.
\end{align*}
代入近轴方程，剩余$(1-q')[2\mathrm ik/q-k^2r^2/q^2]=0$，要求$q'=1$。以束腰处$q(z_w)=-\mathrm iz_R$为初值，$q(z)-q(z_w)=\int_{z_w}^z\dd s=z-z_w$，故
\[
 q=z-z_w-\mathrm iz_R,\qquad z_R>0,
 \quad\frac1q=\frac1{R_w}+\mathrm i\frac2{kw^2}.
\]
这里$z_w$为束腰位置，$w$是场振幅降至轴上$1/\ee$的半径，也即强度降至$1/\ee^2$的半径；$R_w$为波前有向曲率半径。这个$q$的虚部为负，与所选传播相位一致，不能与另一套相位约定的$q=z+\mathrm iz_R$混用。

由复数倒数直接比较实虚部，得到
\[
 w^2(z)=w_0^2\left[1+\frac{(z-z_w)^2}{z_R^2}\right],
 \quad z_R=\frac{kw_0^2}{2}=\frac{\pi n w_0^2}{\lambda_0},
\]
\[
 R_w=(z-z_w)+\frac{z_R^2}{z-z_w}.
\]
束腰处$R_w$无限大，只表示该处波前近乎平面，并不意味着光束以后不发散。远处的发散半角为$w_0/z_R=\lambda_0/(\pi n w_0)$。
\ExpansionDiagram{fig:dep-o1-gaussian}{束腰附近的光束边界；虚线为远区渐近线}{
\draw[->](-3.4,0)--(3.5,0)node[right]{$z-z_w$};\draw[thick,domain=-3:3,samples=90]plot(\x,{.55*sqrt(1+(\x)^2)});\draw[thick,domain=-3:3,samples=90]plot(\x,{-.55*sqrt(1+(\x)^2)});\draw[dashed](-3,1.65)--(3,-1.65);\draw[dashed](-3,-1.65)--(3,1.65);\draw[<->](0,-.55)--(0,.55);\node[right,fill=white,inner sep=1pt]at(.12,.23){$2w_0$};}
图\ref{fig:dep-o1-gaussian}的上下曲线分别为$\pm w(z)$，不是硬边界；边界外仍有指数尾部。计算遮挡功率时要积分实际强度，不能把$\pi w^2$内外简单当作有光与无光。

\subsection{同一ABCD矩阵怎样作用于复参数}
在同一均匀介质的光线坐标$(y,\theta)$下，自由传播使$q\mapsto q+d$。薄透镜附加相位$\exp[-\mathrm ikr^2/(2f)]$，使$1/q\mapsto1/q-1/f$。两者与光线矩阵一致，因此由这些元件组成的系统满足
\[
 q_{\rm out}=\frac{Aq_{\rm in}+B}{Cq_{\rm in}+D}.
\]
可逐个乘矩阵验证分式变换的复合规律。若使用约化坐标$(y,n\theta)$或跨越不同外部介质，应同步改变复参数定义，不能在矩阵中保留折射率而在$q$中忽略它。本节例题两侧均为空气。

\begin{example}透镜后的束腰未必在几何焦平面
一高斯光束恰在薄正透镜平面有束腰，腰半径$w_0$、Rayleigh长度$z_R$，透镜焦距$f>0$。求透镜后的新束腰位置和腰半径，比较$z_R\gg f$与$z_R=f$两种情况。
\ExpansionDiagram{fig:dep-o1-gaussfocus-q}{透镜位于原束腰面，入射波前平坦但束宽有限}{
\draw[->](-3,0)--(4.5,0)node[right]{$z$};\draw[<->,thick](0,-1.4)--(0,1.4)node[above]{透镜};\draw[thick,domain=-2.5:0,samples=70]plot(\x,{.7*sqrt(1+(\x/3)^2)});\draw[thick,domain=-2.5:0,samples=70]plot(\x,{-.7*sqrt(1+(\x/3)^2)});\draw[dashed](3,-1.2)--(3,1.2)node[above]{$f$};\node[left]at(0,-1.6){原束腰};}
\end{example}
\begin{solution}
图\ref{fig:dep-o1-gaussfocus-q}中入射$q_-=-\mathrm iz_R$。薄透镜后
\[
 \frac1{q_+}=-\frac1f+\frac{\mathrm i}{z_R},\qquad
 q_+=-\frac{fz_R^2}{f^2+z_R^2}
 -\mathrm i\frac{f^2z_R}{f^2+z_R^2}.
\]
继续传播距离$z$后，$q(z)=q_++z$。新束腰要求实部为零，所以
\[
 z_w'=\frac{fz_R^2}{f^2+z_R^2},\qquad
 z_R'=\frac{f^2z_R}{f^2+z_R^2},\qquad
 \frac{w_0'}{w_0}=\sqrt{\frac{z_R'}{z_R}}
 =\frac f{\sqrt{f^2+z_R^2}}.
\]
\ExpansionDiagram{fig:dep-o1-gaussfocus-a}{新束腰在几何焦点之前；图示$z_R=f$时位于$f/2$}{
\draw[->](0,0)--(6,0)node[right]{$z$};\draw[thick,domain=0:5,samples=90]plot(\x,{.55*sqrt(1+(\x-2)^2/4)});\draw[thick,domain=0:5,samples=90]plot(\x,{-.55*sqrt(1+(\x-2)^2/4)});\draw[<->](0,-1.5)--(0,1.5);\draw[dashed](2,-1.2)--(2,1.2)node[above]{$f/2$};\draw[dashed](4,-1.5)--(4,1.5)node[above]{$f$};}
当$z_R\gg f$，$z_w'\simeq f$、$w_0'\simeq\lambda_0f/(\pi w_0)$，几何焦平面近似成立。若$z_R=f$，图\ref{fig:dep-o1-gaussfocus-a}中$z_w'=f/2,z_R'=f/2,w_0'=w_0/\sqrt2$。原波前虽平坦，有限腰半径意味着角谱本来就有宽度，不能把它等同于所有光线严格平行的无限平面波。

该结果假定透镜孔径足够大、无明显像差，且新束腰仍满足近轴条件$\lambda_0/(\pi w_0')\ll1$。若透镜把光束截断，单一高斯$q$参数已不能表示完整衍射场，应回到带实际孔径的积分。
\end{solution}

\section{习题与解答}
\begin{exercise}\label{ex:o1oval}
要用一个折射界面把介质$n_1$中的点$A$严格成像到介质$n_2$中的点$B$。忽略反射，证明等光程曲面$n_1|Q-A|+n_2|B-Q|=C$的适当光滑部分满足折射定律，说明与球面近轴成像的区别。
\Diagram{fig:o1ovalq}{笛卡尔卵形面的一个分支，图取$n_1=1,n_2=1.5,C=7.5$}{\fill(0,0)circle(1.5pt)node[below]{$A$};\fill(5.5,0)circle(1.5pt)node[below]{$B$};\draw[curve]plot coordinates {(3.732186,-1.500000) (3.532345,-1.462025) (3.375574,-1.424051) (3.242276,-1.386076) (3.124437,-1.348101) (3.017837,-1.310127) (2.919932,-1.272152) (2.829045,-1.234177) (2.744007,-1.196203) (2.663963,-1.158228) (2.588269,-1.120253) (2.516428,-1.082278) (2.448047,-1.044304) (2.382811,-1.006329) (2.320466,-0.968354) (2.260808,-0.930380) (2.203667,-0.892405) (2.148908,-0.854430) (2.096420,-0.816456) (2.046116,-0.778481) (1.997930,-0.740506) (1.951812,-0.702532) (1.907731,-0.664557) (1.865669,-0.626582) (1.825624,-0.588608) (1.787608,-0.550633) (1.751645,-0.512658) (1.717772,-0.474684) (1.686040,-0.436709) (1.656510,-0.398734) (1.629252,-0.360759) (1.604347,-0.322785) (1.581883,-0.284810) (1.561951,-0.246835) (1.544645,-0.208861) (1.530055,-0.170886) (1.518266,-0.132911) (1.509353,-0.094937) (1.503375,-0.056962) (1.500375,-0.018987) (1.500375,0.018987) (1.503375,0.056962) (1.509353,0.094937) (1.518266,0.132911) (1.530055,0.170886) (1.544645,0.208861) (1.561951,0.246835) (1.581883,0.284810) (1.604347,0.322785) (1.629252,0.360759) (1.656510,0.398734) (1.686040,0.436709) (1.717772,0.474684) (1.751645,0.512658) (1.787608,0.550633) (1.825624,0.588608) (1.865669,0.626582) (1.907731,0.664557) (1.951812,0.702532) (1.997930,0.740506) (2.046116,0.778481) (2.096420,0.816456) (2.148908,0.854430) (2.203667,0.892405) (2.260808,0.930380) (2.320466,0.968354) (2.382811,1.006329) (2.448047,1.044304) (2.516428,1.082278) (2.588269,1.120253) (2.663963,1.158228) (2.744007,1.196203) (2.829045,1.234177) (2.919932,1.272152) (3.017837,1.310127) (3.124437,1.348101) (3.242276,1.386076) (3.375574,1.424051) (3.532345,1.462025) (3.732186,1.500000)};\draw[curve](0,0)--(2.372225,1)--(5.5,0);\draw[guide](0,0)--(2.372225,-1)--(5.5,0);\node[above]at(2.372225,1){$Q$};}
\end{exercise}
\begin{solution}
定义$F(Q)=n_1r_1+n_2r_2$，沿入射传播方向$\vct t_1=(Q-A)/r_1$、出射方向$\vct t_2=(B-Q)/r_2$。梯度为$\nabla F=n_1\vct t_1-n_2\vct t_2$。图\ref{fig:o1ovalq}的等值面法线沿$\nabla F$，任何切向位移$\delta\vct q$满足$\nabla F\cdot\delta\vct q=0$，所以$n_1\vct t_{1\parallel}=n_2\vct t_{2\parallel}$。
\Diagram{fig:o1ovala}{等光程面的法线由入射、出射约化方向矢量之差确定}{\draw[axis](0,0)--(2.4,1.2)node[above left]{$n_1\vct t_1$};\draw[axis](0,0)--(2.1,-.5)node[below]{$n_2\vct t_2$};\draw[->,thick](2.1,-.5)--(2.4,1.2);\node[right]at(2.6,.35){差矢量};\draw[guide](0,0)--(.3,1.7);}
上式为斯涅尔切向条件，法线的矢量构造见图\ref{fig:o1ovala}。须选择使光由介质1进入介质2、无全反射且不被其他界面截断的分支。球面一般只能在轴上小孔径处近似满足固定物像等光程；这个等光程面通常不是球面，可以为指定物像对校正球差，却不保证别的物点也严格成像。
\end{solution}

\begin{exercise}\label{ex:o1galilean}
空气中正物镜焦距$f_1>0$，负目镜$f_2<0$，间隔$d=f_1+f_2>0$。求无焦系统角放大率、轴向平行束宽度比和像的正倒，与双正透镜系统比较。
\Diagram{fig:o1galileoq}{负目镜截取尚未到物镜焦点的会聚光束}{\OptAxis\draw[<->](1,-1.3)--(1,1.3);\draw[>-<](4,-1)--(4,1);\draw[guide](5.5,-1.5)--(5.5,1.5);\node[above]at(5.5,1.5){物镜焦点};\node[above]at(4,1){负目镜};}
\end{exercise}
\begin{solution}
图\ref{fig:o1galileoq}仍用两透镜矩阵，把$d=f_1+f_2$代入得到
\[
 M=\begin{pmatrix}-f_2/f_1&f_1+f_2\\0&-f_1/f_2\end{pmatrix}.
\]
因$f_2<0$，两个对角元都正。轴向平行光的宽度比$|f_2|/f_1$，小视角放大率$\theta_o/\theta_i=f_1/|f_2|>0$，所以角方向不反转。
\Diagram{fig:o1galileoa}{光束在到达几何焦点之前被负透镜重新准直}{\OptAxis\draw[<->](1,-1.3)--(1,1.3);\draw[>-<](4,-1)--(4,1);\draw[->](0,1)--(1,1)--(4,.333333)--(6,.333333);\draw[->](0,-1)--(1,-1)--(4,-.333333)--(6,-.333333);\draw[guide](4,.333333)--(5.5,0);\draw[guide](4,-.333333)--(5.5,0);}
图\ref{fig:o1galileoa}按$f_1=4.5,f_2=-1.5,d=3$作图，束宽缩为1/3、角放大率3，内部没有经过实焦点。双正系统角放大率为负，会产生倒向角像。两系统都满足$\Phi=0$，但角放大率符号和内部光路不同，对应正立、倒立两种观察效果。
\end{solution}

\begin{exercise}\label{ex:o1layers}
平行分层透明介质的各层厚度为$h_j$、折射率$n_j$。底部点物体从顶面外的介质$n_0$中近轴观察。求表观深度，说明为何各层的真实光程厚度不能直接作为表观深度相加。
\Diagram{fig:o1layersq}{平行分层结构保持每条光线的切向约化分量}{\foreach\y in{0,.7,1.4,2.1}{\draw[curve](-2,\y)--(2,\y);}\fill(0,0)circle(1.5pt)node[below]{物点};\draw[->](0,0)--(.35,.7)--(.6,1.4)--(1.1,2.1)--(1.8,3);\node[right]at(2,.35){$n_3,h_3$};\node[right]at(2,1.05){$n_2,h_2$};\node[right]at(2,1.75){$n_1,h_1$};}
\end{exercise}
\begin{solution}
图\ref{fig:o1layersq}所有平面界面无光焦度，$p=n\theta$在近轴下保持。光从底部穿过第$j$层的横向位移为$h_j\theta_j=h_jp/n_j$，到顶面总横移$y=p\sum_jh_j/n_j$。外部出射斜率为$\theta_0=p/n_0$。
\Diagram{fig:o1layersa}{外部光线反向延长至虚像，深度由横移除外部斜率决定}{\draw[curve](-2,1.6)--(2,1.6);\fill(0,0)circle(1.5pt)node[below]{真实物};\fill(0,.6)circle(1.5pt)node[left]{虚像};\draw[guide](0,.6)--(1,1.6);\draw[->](1,1.6)--(2.3,2.9);\draw[guide](0,0)--(0,2.6);}
图\ref{fig:o1layersa}表观深度$h_a=y/\theta_0=n_0\sum_jh_j/n_j$。真实传播光程为$\sum_jn_jh_j$，一个含$n_j$倒数、另一个含$n_j$本身，物理意义不同。若各层折射率相同$n$，便恢复$n_0h/n$；若$n_0=n$，深度不变。
\end{solution}

\begin{exercise}\label{ex:o1immersion}
同一薄双凸透镜折射率$n_l=1.50$，在空气中焦距$f_a$。完全浸入折射率$n_0=1.33$的液体，求主面焦距与$f_a$之比，讨论液体折射率超过玻璃时的光焦度符号。
\Diagram{fig:o1immersionq}{同一曲率透镜置于不同外部介质}{\OptAxis\draw[curve](2,-1.3)..controls(1.4,-.4)and(1.4,.4)..(2,1.3);\draw[curve](2,-1.3)..controls(2.6,-.4)and(2.6,.4)..(2,1.3);\node[above]at(2,1.5){$n_l$};\node at(4,1){$n_0$};}
\end{exercise}
\begin{solution}
令曲率组合$K=1/R_1-1/R_2>0$。图\ref{fig:o1immersionq}在外部介质中约化光焦度$\Phi=(n_l-n_0)K$，两侧相同折射率的主面焦距$f_0=n_0/\Phi$。空气中$f_a=1/[(n_l-1)K]$，所以
\[
 \frac{f_0}{f_a}=\frac{n_0(n_l-1)}{n_l-n_0}
 =\frac{1.33\times.50}{.17}\simeq3.912.
\]
\Diagram{fig:o1immersiona}{外介质接近透镜折射率，折射偏折变弱、焦距增大}{\OptAxis\draw[<->](1.5,-1.2)--(1.5,1.2);\draw[->](0,.8)--(1.5,.8)--(5.8,0);\draw[guide](1.5,.8)--(2.6,0);\node[below]at(2.6,0){空气焦点};\node[below]at(5.8,0){液体焦点};}
图\ref{fig:o1immersiona}焦距增大是相对折射率差减小的结果。$n_0=n_l$时界面不可光学分辨，薄透镜光焦度为零；$n_0>n_l$时$\Phi<0$，相同双凸形状转为发散。透镜的会聚、发散性质由相对折射率和几何形状共同决定。
\end{solution}

\begin{exercise}\label{ex:o1mirage}
平移不变介质中$n(y)=n_A-gy$，$g>0$，只取$n>0$区域。光从$(z,y)=(0,0)$向上方以相对$+z$角$\alpha_0\in(0,\pi/2)$进入。求精确轨迹、最高点和回到$y=0$的水平距离。
\Diagram{fig:o1mirageq}{折射率向上减小，光线弯回高折射率区}{\OptAxis\draw[axis](0,0)--(0,3)node[above]{$y$};\draw[axis](0,0)--(1.4,1)node[above]{$\alpha_0$};\draw[guide](0,2.5)--(5.5,2.5);\node[right]at(5.5,2.5){$n$较小};}
\end{exercise}
\begin{solution}
图\ref{fig:o1mirageq}的守恒量$b=n\cos\alpha=n_A\cos\alpha_0$。轨迹斜率$y'=\tan\alpha$，所以$n^2=b^2(1+y'^2)$。在上升支$y'>0$，令$n=b\cosh u$，则$y'=\sinh u$；由$n=n_A-gy$，$\dd y=-(b/g)\sinh u\dd u$，于是$\dd z=-(b/g)\dd u$。
\Diagram{fig:o1miragea}{轨迹是倒置悬链线的一段，对最高点左右对称}{\OptAxis\draw[axis](0,0)--(0,2.7)node[above]{$y$};\draw[curve,domain=0:2.8,samples=80]plot(\x,{1.4*(cosh(1)-cosh((\x-1.4)/1.4))});\draw[guide](1.4,0)--(1.4,.760313);\node[above]at(1.4,.85){最高点};\node[below]at(2.8,0){返回};}
以最高点$u=0,z=z_t$为参考点，上升支满足
\[
 \int_{z_t}^{z}\dd Z=-\frac bg\int_0^u\dd q,
 \qquad z-z_t=-\frac bg u.
\]
故$u=g(z_t-z)/b$。由初始条件$n(0)=n_A=b\cosh(gz_t/b)$求出$z_t$，再以$\cosh$的偶性接续下降支，得到
\[
 y(z)=\frac{n_A-b\cosh[g(z-z_t)/b]}g,\quad
 z_t=\frac bg\operatorname{arcosh}\frac{n_A}b.
\]
最高点$y_t=(n_A-b)/g$，回到$y=0$时$z=2z_t$。图\ref{fig:o1miragea}按一组比例画出同一悬链线结构。这个解适用于有限角度。前面的$n^2$抛物分布也有精确的正弦型轨迹，但其轴向周期依赖初始方向；仅在近轴极限下才化为固定节距$2\pi a$。两种分布的轨迹都须限制在实际介质覆盖的区域内。
\end{solution}

\begin{exercise}\label{ex:o1chromatic}
薄透镜几何曲率固定，空气中折射率随波长变化。已知$n(\lambda_0)=n_0$，一小段波长变化对应$\Delta n$。求焦距一阶变化及平行白光的轴向色差方向。
\Diagram{fig:o1chromq}{同一透镜对不同波长的光焦度不同}{\OptAxis\draw[<->](1.8,-1.3)--(1.8,1.3);\draw[->](0,1)--(1.8,1);\draw[curve](1.8,1)--(4,0);\draw[guide](1.8,1)--(5.6,0);\node[below]at(4,0){短波};\node[below]at(5.6,0){长波};}
\end{exercise}
\begin{solution}
图\ref{fig:o1chromq}的$f=K^{-1}(n-1)^{-1}$。固定$K$求导，$\dd f/\dd n=-K^{-1}(n-1)^{-2}=-f/(n-1)$，所以$\Delta f\simeq-f_0\Delta n/(n_0-1)$。对题图中的正透镜，正常色散使短波折射率较大，蓝光因而比红光更早会聚；发散透镜也服从同一有向焦距微分式。
\Diagram{fig:o1chroma}{像面固定时，不同焦点在屏上形成不同半径模糊斑}{\OptAxis\draw[<->](1,-1.2)--(1,1.2);\draw[guide](4,-1.6)--(4,1.6);\draw[curve](1,.9)--(3,0)--(5,-.9);\draw[curve](1,-.9)--(3,0)--(5,.9);\draw[guide](1,.9)--(4.5,0);\draw[guide](1,-.9)--(4.5,0);}
图\ref{fig:o1chroma}说明固定屏幕不可能同时处在所有波长的几何焦面。色差是$n(\lambda)$改变光焦度的几何效应，单色衍射则即使焦距相同仍形成有限光斑；两种宽度不能不分机制地直接相加。本一阶式还要求$|\Delta n|\ll n_0-1$。
\end{solution}

\begin{exercise}\label{ex:o1prismminimum}
折射率$n$、顶角$A$的棱镜在空气中，光由一个斜面入射、另一个出射。推导对称路径下的最小偏向角公式，说明该状态需满足的传播条件。
\Diagram{fig:o1prismq}{一般穿透光路；法线与界面垂直，入射、折射角分别对法线量取}{
\draw[curve](0,0)--(2,3.4641)--(4,0)--cycle;
\draw[->](-.7,.395855)--(.692820,1.2);
\draw[curve](.692820,1.2)--(3.175128,1.428721);
\draw[->](3.175128,1.428721)--(4.7,1.190565);
\draw[guide](0,1.6)--(1.385641,.8);
\draw[guide](2.482307,1.028721)--(3.867948,1.828721);
\node at(2,2.75){$A$};\node at(2,.55){$n>1$};
\node[above left]at(.28,1.23){$i_1$};\node[below right]at(.98,1.2){$r_1$};
}
\end{exercise}
\begin{solution}
图\ref{fig:o1prismq}中折射角满足$r_1+r_2=A$，空气侧入、出射角满足$\sin i_j=n\sin r_j$，总偏向为$D=i_1+i_2-A$。令$f(r)=\arcsin(n\sin r)$，则$D=f(r_1)+f(A-r_1)-A$。在$n>1$、$0\leq r_j<\arcsin(1/n)$的通常穿透分支，
\begin{align*}
 f'(r)&=\frac{n\cos r}{\sqrt{1-n^2\sin^2r}},\\
 f''(r)&=-\frac{n\sin r}{\sqrt{1-n^2\sin^2r}}
 +\frac{n^3\sin r\cos^2r}{(1-n^2\sin^2r)^{3/2}}\\
 &=\frac{n(n^2-1)\sin r}{(1-n^2\sin^2r)^{3/2}}\geq0.
\end{align*}
凸性保证$r_1=r_2=A/2$为极小点，光路可逆性则给出对称驻点。
\Diagram{fig:o1prisma}{最小偏向的对称光路；图取顶角60度、折射率1.5}{
\draw[curve](0,0)--(2,3.4641)--(4,0)--cycle;
\draw[->](-.7,.8929)--(.8083,1.4);
\draw[->](.8083,1.4)--(3.1917,1.4);
\draw[->](3.1917,1.4)--(4.7,.8929);
\draw[guide](2,-.2)--(2,3.7);
\draw[guide](.1155,1.8)--(1.5011,1);
\draw[guide](2.4989,1)--(3.8845,1.8);
\node at(1.3,1.1){$r_1$};\node at(2.7,1.1){$r_2$};
\node[below]at(2,-.3){$r_1=r_2=A/2$};
\node at(2,2.85){$A$};
}
图\ref{fig:o1prisma}的对称路径有$i_1=i_2=i$，因此
\[
 D_{\min}=2i-A,
 \qquad n=\frac{\sin[(A+D_{\min})/2]}{\sin(A/2)}.
\]
可传播条件为$n\sin(A/2)\leq1$；等号对应掠入、掠出极限，有限通光通常要求严格小于1。若不满足，无法让光沿这个两面穿透分支通过。薄棱镜$A\ll1$且$nA\ll1$时$i\simeq nA/2$，于是$D_{\min}\simeq(n-1)A$。近似要求两个空气侧角度也小，不能只检查棱镜顶角而忽略很大的折射率。
\end{solution}

\begin{exercise}\label{ex:o1obliqueplate}
平面波以角$i$穿过厚$h$、折射率$n$的平行板，板外为空气。比较板后同一平面上的相位与未插板情况，推导附加相位，避免把折线实际路长误当程差。
\Diagram{fig:o1obliqueplateq}{比较同一输出平面上的波相位}{\draw[curve](-2,0)--(2,0);\draw[curve](-2,1)--(2,1);\draw[->](-1.5,-1)--(-.5,0)--(.1,1)--(1.1,2);\draw[guide](-.5,0)--(.5,1);\node[right]at(2,.5){$h$};}
\end{exercise}
\begin{solution}
图\ref{fig:o1obliqueplateq}中平行界面保持切向波矢，$\sin i=n\sin r$。在固定横坐标比较上下参考面的相位增长，板内法向波矢$k_{zg}=k_0n\cos r$，原空气中为$k_{za}=k_0\cos i$。所以附加相位$\Delta\phi=k_0h(n\cos r-\cos i)$。
\Diagram{fig:o1obliqueplatea}{沿不同几何射线积累的光程须先归到同一参考波面}{\draw[curve](0,0)--(4,0);\draw[curve](0,1)--(4,1);\draw[guide](1,-.5)--(1,1.7);\draw[->](.3,-.6)--(1,0)--(1.5,1)--(2.2,1.6);\node[above]at(1,1.8){固定横坐标};}
图\ref{fig:o1obliqueplatea}的计算自动扣除了横向位移相关相位。正常入射恢复$k_0(n-1)h$。若题目改为固定两个点之间的封闭干涉臂，需同时满足光线连接、元件位置和合束条件，不能只把本平面波相位式套给任意折线路径。
\end{solution}

\begin{exercise}\label{ex:o1matrixinverse}
空气近轴系统矩阵为$M$且$\det M=1$。写出逆矩阵，解释“数学逆传播”与把整套透镜实际前后翻转为何须先统一坐标基底。
\Diagram{fig:o1inverseq}{同一对参考面之间的前向、反向关系}{\draw[curve](1,-1)rectangle(3.5,1);\node at(2.25,0){$M$};\draw[->](-.5,.5)--(.8,.5);\draw[->](3.7,-.5)--(5,-.5);}
\end{exercise}
\begin{solution}
图\ref{fig:o1inverseq}用同一固定$z$和有向倾角描述时，数学逆为$M^{-1}=\bigl(\begin{smallmatrix}D&-B\\-C&A\end{smallmatrix}\bigr)$，因为与原矩阵相乘给单位阵。逆关系把末端的$y,n\theta$恢复到起点，并不自动改变倾角的参考轴。
\Diagram{fig:o1inversea}{反向实际传播若重新令光轴向前，倾角符号也改变}{\draw[axis](0,.7)--(4,.7)node[right]{$+z$};\draw[axis](4,-.7)--(0,-.7)node[left]{$+\tilde z$};\node at(2,0){同一光线需换基};}
若反向传播时重新把轴指向左方而高度仍向上，约化倾角需变号，换基矩阵$K=\operatorname{diag}(1,-1)$，反向系统矩阵为$KM^{-1}K=\bigl(\begin{smallmatrix}D&B\\C&A\end{smallmatrix}\bigr)$。图\ref{fig:o1inversea}说明相同光学系统反向后的$B,C$正负取决于坐标约定，不能只看一处符号判定互易性失效。
\end{solution}

\begin{exercise}\label{ex:o1focusplate}
物点到固定屏距离$D$，中间放焦距$f$的薄透镜，空气中近轴成像。求允许的透镜位置数及放大率乘积，讨论$D<4f$。
\Diagram{fig:o1focusq}{固定物屏间距，透镜位置待定}{\OptAxis\draw[->](0,0)--(0,1);\draw[<->](2,-1.2)--(2,1.2);\draw[curve](6,-1.5)--(6,1.5);\draw[<->](0,-1.8)--(6,-1.8)node[midway,below]{$D$};}
\end{exercise}
\begin{solution}
图\ref{fig:o1focusq}中取物体在0、屏在$D$，透镜位置为$u$，于是像距$v=D-u$。由$1/u+1/v=1/f$，相乘得到$uv=f(u+v)=fD$，即
\[
 u^2-Du+fD=0,\qquad
 u_\pm=\frac{D\pm\sqrt{D^2-4fD}}2.
\]
题中的实物、实屏位于透镜两侧，故要求$0<u<D$；当$D>4f$、$f>0$时两个根均满足此条件。
\Diagram{fig:o1focusa}{两次成像透镜位置关于物屏中点对称}{\OptAxis\draw[<->](1.7,-1.2)--(1.7,1.2);\draw[guide](4.3,-1.2)--(4.3,1.2);\draw[curve](6,-1.5)--(6,1.5);\draw[guide](3,-1.7)--(3,1.7);\node[above]at(1.7,1.2){位置一};\node[above]at(4.3,1.2){位置二};}
图\ref{fig:o1focusa}的两个位置将$u,v$互换，放大率分别为$-v/u$和$-u/v$，乘积1。一幅放大像与一幅缩小像配对出现。两位置间距
\[
 l=u_+-u_- =\sqrt{D^2-4fD},\qquad
 f=\frac{D^2-l^2}{4D}.
\]
这给出了用固定物屏间距与两次清晰成像位置测焦距的方法。$D=4f$时判别式为零，唯一位置$u=v=2f$、放大率$-1$；$D<4f$时没有满足实物实像要求的透镜位置。也可由$u+v\geq2\sqrt{uv}=2\sqrt{fD}$直接得到$D\geq4f$。实际厚透镜要以主面测距，若机械测的是两顶点位置，须先修正主面偏移。
\end{solution}

\begin{exercise}\label{ex:o1etendue}
理想近轴成像系统横向放大率为$m_y$，物像介质折射率分别$n_O,n_I$。求同一轴上物点光束的角放大率，并解释放大像时为何常需减小像方数值孔径。
\Diagram{fig:o1etendueq}{同一物点的两条光线定义孔径角}{\draw[axis](-.3,0)--(6,0);\draw[curve](2,-1.1)rectangle(3,1.1);\draw[curve](0,0)--(2,1);\draw[curve](0,0)--(2,-1);\draw[curve](3,.6)--(5,0);\draw[curve](3,-.6)--(5,0);}
\end{exercise}
\begin{solution}
图\ref{fig:o1etendueq}中，物像面之间使用约化变量$(y,p=n\theta)$。成像要求$B=0$，横向放大率为$A=m_y$。矩阵行列式恒为1，故$AD=1$、$D=1/m_y$，有
\[
 \begin{pmatrix}y_I\\n_I\theta_I\end{pmatrix}
 =\begin{pmatrix}m_y&0\\ C&1/m_y\end{pmatrix}
 \begin{pmatrix}y_O\\n_O\theta_O\end{pmatrix}.
\]
对轴上同一物点$y_O=0$，下行式化为$n_I\theta_I=n_O\theta_O/m_y$，所以角放大率为$n_O/(n_Im_y)$。
\Diagram{fig:o1etenduea}{像高和约化孔径角的乘积保持}{\draw[curve](0,0)rectangle(1.8,.7);\node at(.9,.35){$h_O$};\draw[curve](3,0)rectangle(5.7,1.1);\node at(4.35,.55){$h_I=m_yh_O$};\node[below]at(3,-.8){$n_Oh_O\theta_O=n_Ih_I\theta_I$};}
图\ref{fig:o1etenduea}所画两束使用同一组通光光线，其数值孔径满足$\mathrm{NA}_I\simeq\mathrm{NA}_O/|m_y|$。负放大率还带来倾角反号，孔径大小则取绝对值。对轴上附近小面元，面积放大$|m_y|^2$，小圆锥立体角约为$\pi\theta^2$，于是
\[
 n_I^2\dd A_I\dd\Omega_I
 =n_O^2\dd A_O\dd\Omega_O.
\]
该关系给出近轴通光相空间体积的不变性。无损时同一光束的功率也守恒，因此辐亮度除以$n^2$保持。扩大像面积会减小相应像方角锥；若另行增大物方收集孔径，所比较的已是另一组光线，不能把两种条件下的亮度混为一谈。
\end{solution}

\begin{exercise}\label{ex:o1magnification}
薄透镜在空气中$f>0$，物高$h$，物距$u>f$。求像距与横向放大率，并由两条光线画出$u=5f/3$的像。
\Diagram{fig:o1magnifyq}{已知物高、物距与焦距}{\OptAxis\draw[<->](2,-1.8)--(2,1.8);\draw[->,thick](0,0)--(0,1);\node[left]at(0,1){$h$};\fill(3.2,0)circle(1.3pt)node[below]{$F'$};\draw[<->](0,-.6)--(2,-.6)node[midway,below]{$u$};}
\end{exercise}
\begin{solution}
图\ref{fig:o1magnifyq}中$1/u+1/v=1/f$，解得$v=fu/(u-f)$。通过光心的近轴光线保持方向，由相似三角形$h'/v=-h/u$，所以横向放大率$m_y=h'/h=-v/u=-f/(u-f)$。
\Diagram{fig:o1magnifya}{平行轴光线经后焦点，过光心光线不偏折，两者交于倒立实像}{\OptAxis\draw[<->](2,-1.9)--(2,1.9);\draw[->,thick](0,0)--(0,1);\draw[->,thick](5,0)--(5,-1.5);\draw[->](0,1)--(2,1);\draw[->](2,1)--(5,-1.5);\draw[->](0,1)--(5,-1.5);\fill(3.2,0)circle(1.2pt)node[above]{$F'$};}
取$u=5f/3$，则$v=5f/2$、$m_y=-3/2$。图\ref{fig:o1magnifya}按$f=1.2$、$u=2$作图，实像在$z=5$、像高$-1.5h$，交点与计算吻合。真实透镜在大孔径或大物高时需考虑像差，本图只表达近轴模型。
\end{solution}

\begin{exercise}\label{ex:o1afocal}
正透镜$f_1>0$、$f_2>0$相隔$f_1+f_2$。求平行细光束通过后的束径比与小角偏转比，并说明无焦系统是否“没有作用”。
\Diagram{fig:o1afocalq}{共焦双正透镜系统}{\OptAxis\draw[<->](1,-1.4)--(1,1.4);\draw[<->](5,-1.1)--(5,1.1);\draw[guide](3.5,-1.6)--(3.5,1.6);\node[above]at(3.5,1.6){$F'_1=F_2$};}
\end{exercise}
\begin{solution}
图\ref{fig:o1afocalq}中$d=f_1+f_2$，直接乘矩阵得到
\[
 M=\begin{pmatrix}-f_2/f_1&f_1+f_2\\0&-f_1/f_2\end{pmatrix}.
\]
对于轴向平行光$\theta_i=0$，$y_o=-(f_2/f_1)y_i$，输出仍平行，束径绝对比$f_2/f_1$。
\Diagram{fig:o1afocala}{轴向光束先会聚再重新平行，图取两焦距不同}{\OptAxis\draw[<->](1,-1.4)--(1,1.4);\draw[<->](5,-1.4)--(5,1.4);\draw[->](0,1)--(1,1);\draw[->](1,1)--(5,-.6);\draw[->](5,-.6)--(6.5,-.6);\draw[->](0,-1)--(1,-1);\draw[->](1,-1)--(5,.6);\draw[->](5,.6)--(6.5,.6);}
对整体倾斜的小角平行光，$\theta_o=-(f_1/f_2)\theta_i$，与高度无关。图\ref{fig:o1afocala}取$f_1/f_2=5/3$，中间焦点位于第一透镜后2.5图示单位。无焦仅表示$C=0$、平行束仍为平行束；系统仍可扩束、缩束和改变角度，绝非恒等映射。
\end{solution}

\begin{exercise}\label{ex:o1glassplate}
一块厚度$d$、折射率$n$的平行玻璃板插入原本将于空气中会聚的近轴光束，板面垂直轴。求焦点沿轴的位移，忽略反射。
\Diagram{fig:o1glassq}{会聚光束在平行板内传播的斜率减小}{\OptAxis\draw[curve](2,-1.1)rectangle(3,1.1);\draw[guide](0,1.3)--(4.8,0);\draw[guide](0,-1.3)--(4.8,0);\node[above]at(2.5,1.1){$n,d$};\fill(4.8,0)circle(1.3pt)node[below]{$F_0$};}
\end{exercise}
\begin{solution}
图\ref{fig:o1glassq}中原光线以斜率$\theta<0$会聚。进入平面板后$\theta_g=\theta/n$，传播$d$导致高度变化$d\theta/n$；若仍为空气则变化$d\theta$。出板后倾角恢复$\theta$，故高度比原光线高出$\Delta y=d\theta(1/n-1)>0$。
\Diagram{fig:o1glassa}{出板后两光线平行，新的交轴点向后移动}{\OptAxis\draw[curve](2,-1.1)rectangle(3,1.1);\draw[->](0,1.2)--(2,.7)--(3,.575)--(5.3,0);\draw[guide](2,.7)--(4.8,0);\fill(4.8,0)circle(1.2pt)node[below]{$F_0$};\fill(5.3,0)circle(1.2pt)node[above right]{$F$};}
图\ref{fig:o1glassa}中的轴向位移由$\theta\Delta z=-\Delta y$得$\Delta z=d(1-1/n)>0$。矩阵角度看，是把空气传播长度$d$的约化距离$d$替换为$d/n$，而折射后的外部倾角不变。大孔径时不同光线位移并不完全相同，会引入球差。
\end{solution}

\begin{exercise}\label{ex:o1turning}
水平分层介质$n(y)=n_0\sqrt{1-y^2/a^2}$。光线从$y=0$以与$+z$轴成$\alpha_0$的方向进入。用守恒量求最大高度，不使用近轴条件。
\Diagram{fig:o1turnq}{光线在折射率下降区弯折至水平}{\OptAxis\draw[axis](0,0)--(0,2.7)node[above]{$y$};\draw[curve](0,0)..controls(1,1)and(2,1.8)..(3.2,1.8);\draw[guide](0,1.8)--(4,1.8);\node[left]at(0,1.8){$y_t$};\node at(.8,.3){$\alpha_0$};}
\end{exercise}
\begin{solution}
图\ref{fig:o1turnq}中，介质沿$z$平移不变，$b=n\cos\alpha=n_0\cos\alpha_0$守恒。用$y'=\tan\alpha$表示方向，得到
\[
 n^2(y)=b^2(1+y'^2),\qquad
 y'^2=\frac{\sin^2\alpha_0-y^2/a^2}{\cos^2\alpha_0}.
\]
真实射线要求右侧非负，允许区域为$|y|\leq a|\sin\alpha_0|$。最高点$y'=0$，所以正向上射线的最大高度为$y_t=a\sin\alpha_0$；若只讨论幅度，则为$a|\sin\alpha_0|$。
\Diagram{fig:o1turna}{允许高度由$n(y)\geq n_0\cos\alpha_0$确定}{\draw[axis](0,0)--(4.2,0)node[right]{$y/a$};\draw[axis](0,0)--(0,2.8)node[above]{$n/n_0$};\draw[curve,domain=0:.99,samples=70]plot({3.5*\x},{2.3*sqrt(1-\x*\x)});\draw[guide](0,1.7)--(3.7,1.7);\node[left]at(0,1.7){$\cos\alpha_0$};}
图\ref{fig:o1turna}中的折射率曲线与水平$b$线相交处就是转点。还可将方程积分，以核对近轴极限。对$0<\alpha_0<\pi/2$的上升支，令$A=a\sin\alpha_0$，则
\[
 \int_0^{y}\frac{\dd q}{\sqrt{A^2-q^2}}
 =\frac1{a\cos\alpha_0}\int_0^z\dd s,\qquad
 [\arcsin(q/A)]_0^{y}=\frac{z}{a\cos\alpha_0}.
\]
这里用$q=A\sin u$，在上升支$0\leq u\leq\pi/2$有$\dd q=A\cos u\dd u$、$\sqrt{A^2-q^2}=A\cos u$。反解得$y=A\sin[z/(a\cos\alpha_0)]$，在转点处接续下降支。
精确轴向周期是$2\pi a\cos\alpha_0$，最高点在$z_t=\pi a\cos\alpha_0/2$。小角时振幅$A\simeq a\alpha_0$、周期趋$2\pi a$，恢复近轴解。有限角度下转点高度与周期都应使用轴向守恒量求出，不能仅保留振幅修正而仍把近轴周期称为精确周期。
\end{solution}

\begin{exercise}\label{ex:o1centralturn}
圆对称介质在某区域$n(r)=n_a a/r$。求一条非径向光线能否有孤立最近点，推导其轨迹形状。讨论守恒量$b=n_a a$的特殊情形。
\Diagram{fig:o1centralq}{径向折射率场中的光线与径向夹角}{\draw[curve](0,0)circle(1);\draw[curve](0,0)circle(2);\draw[axis](0,0)--(2,1)node[right]{$r$};\draw[axis](2,1)--(2.5,2.2)node[above]{切向};\node[below left]at(0,0){$O$};}
\end{exercise}
\begin{solution}
图\ref{fig:o1centralq}的$n(r)r=n_aa$为常数，所以$\sin\psi=b/(n_aa)$也为常数。若$0<b<n_aa$，夹角固定且不等于直角，不存在孤立转点。由$\dd r/\dd s=\cos\psi$、$r\dd\varphi/\dd s=\sin\psi$，得$\dd\varphi/\dd r=\tan\psi/r$。
\Diagram{fig:o1centrala}{固定夹角对应对数螺线，图只画有限介质区域}{\draw[axis](-2.3,0)--(2.5,0);\draw[axis](0,-2.3)--(0,2.5);\draw[curve,domain=0:300,samples=100]plot({.35*exp(.006*\x)*cos(\x)},{.35*exp(.006*\x)*sin(\x)});\node[below left]at(0,0){$O$};}
以$(r_0,\varphi_0)$为初始点积分，$\int_{\varphi_0}^{\varphi}\dd\phi=\tan\psi\int_{r_0}^{r}\dd q/q=\tan\psi[\ln q]_{r_0}^{r}$，得$\varphi-\varphi_0=\tan\psi\ln(r/r_0)$，即图\ref{fig:o1centrala}的对数螺线。向内支的$\cos\psi<0$，须保留$\tan\psi$的符号。若$b=n_aa$，则$\psi=\pi/2$，光线可保持任意给定半径的圆周轨道；此时不能再除以$\dd r/\dd s$，应直接使用切向关系。该折射率在$r=0$奇异，轨迹仅适用于题定介质区域。
\end{solution}

\begin{exercise}\label{ex:o1phasors}
同频同偏振两束光复振幅为$A$和$2A\ee^{i\delta}$。求合振幅的大小、相位以及最大最小强度，并说明相位求值应怎样处理象限。
\Diagram{fig:o1phasorq}{两相量大小之比为1:2}{\draw[axis](0,0)--(1.5,0)node[right]{$A$};\draw[axis](0,0)--(1.5,2.598076)node[right]{$2A\ee^{i\delta}$};\draw(.6,0)arc(0:60:.6);\node at(.8,.4){$\delta$};}
\end{exercise}
\begin{solution}
以第一束的相位为零，图\ref{fig:o1phasorq}中合复振幅为
\[
 \frac{\mathcal E}A=1+2\ee^{i\delta}
 =1+2\cos\delta+2i\sin\delta.
\]
因此$|\mathcal E|^2/A^2=(1+2\cos\delta)^2+4\sin^2\delta=5+4\cos\delta$，合振幅大小为$A\sqrt{5+4\cos\delta}$，相位为
\[
 \phi=\operatorname{atan2}(2\sin\delta,1+2\cos\delta).
\]
该函数由实、虚两分量共同确定象限；单独取反正切商只能给出相差整数$\pi$的一组可能值。
\Diagram{fig:o1phasora}{合相量实部可能为负，单值反正切会丢失象限}{\begin{scope}[scale=.7]\draw[axis](-2.2,0)--(3.4,0)node[right]{实部};\draw[axis](0,-2.3)--(0,2.8)node[above]{虚部};\draw[curve](1,0)circle(2);\draw[->,thick](0,0)--(-.5,1.322876);\node[left]at(-.5,1.322876){$\mathcal E$};\end{scope}}
图\ref{fig:o1phasora}的合相量末端以$(A,0)$为圆心、$2A$为半径。原点位于圆内，合相量会经过左半平面，但不会为零，因为两束振幅不等。若$A$对应单束强度$I_1$，则
\[
 I_{\max}=9I_1\quad(\delta=2m\pi),\qquad
 I_{\min}=I_1\quad(\delta=(2m+1)\pi),
\]
可见度为$(9-1)/(9+1)=4/5$。合相量越过虚轴、实部变号时，相位连续经过$\pm\pi/2$；越过负实轴时，取值区间为$(-\pi,\pi]$的主值相位才会发生$2\pi$跳变。本题在$\delta=\pi$通过负实轴，可按连续相位消去这种记号上的跳变。场本身始终连续且不为零，相位也始终有定义。
\end{solution}

